\section{Radial parts and the Moyal normalization}
\label{sec:app-radial-parts-moyal}

Formal Weyl quantization of the constant symplectic plane and stable
scalar-reduced trace map on the Chan--Paton brane stack disagree
at each bidegree \((a,b)\) by a single class.  The disagreement lives
in a finite-dimensional bidegree-local complex whose source records
cyclic-word boundary letters and whose target combines cyclic-word
incidence with the residual normal-diagram corrections.  Concretely,
\(\mathsf W_{a,b}\) is the rational span of words containing the
moment-map letter \([x,y]\) once;
\(\mathsf V_{a,b}\) is the cyclic-word target of the necklace
boundary; \(\mathsf{Diag}_{a,b}\) is the free indexed normal-diagram
space carrying contraction-degree-\(\geq1\) (PBW) corrections; and
\(B_{a,b}\) is the combined coboundary
\((\partial,C^+_{a,b})\) of
Definition~\ref{def:app-free-kernel-correction-obstruction}.  The
discrepancy is then
\[
  \Omega^{\mathrm{rad}}_{a,b}
    \in
  \operatorname{coker}B_{a,b},
  \qquad
  B_{a,b}\colon\mathsf W_{a,b}
    \longrightarrow
  \mathsf V_{a,b}\oplus\mathsf{Diag}_{a,b},
\]
in the bidegree-local radial bicomplex
\((\mathsf W_{a,b}\xrightarrow{\partial}\mathsf V_{a,b},\;
   \mathsf W_{a,b}\xrightarrow{C^+_{a,b}}\mathsf{Diag}_{a,b})\)
of Definition~\ref{def:app-free-kernel-correction-obstruction}.  Its
vanishing has one unconditional finite-window form and one conditional
all-bidegree form.  Proposition~\ref{prop:app-radial-pr-unconditional-window}
gives exact rational vanishing on the verified finite window listed in
Table~\ref{tab:app-radial-finite-window-status}.  The all-bidegree
statement is Theorem~\ref{thm:app-radial-bicomplex-acyclic} only under
the radial-parts datum of
Statement~\ref{stmt:app-radial-parts-hypotheses}.  Outside these two
inputs, \(\Omega^{\mathrm{rad}}_{a,b}\) is the obstruction class.

The cyclic-word incidence complex
\(\partial:\mathsf W_{a,b}\to\mathsf V_{a,b}\) is the \(\hbar_{\mathrm{W}}^0\)
skeleton of the radial bicomplex; the classical lift
\(\partial c=T_{a,b}\) is the ordinary necklace Stokes identity of
Lemma~\ref{lem:app-ordinary-necklace-stokes} (\(T_{a,b}\) the necklace
target of a moment-map insertion).  The free indexed
normal-diagram target \(\mathsf{Diag}_{a,b}\) records the residual
contraction-degree-\(\geq1\) (PBW) corrections and the Capelli
\(\hbar_{\mathrm{W}} N\) term; the residual class
\(R^{\mathrm{free}}_{a,b}\in\mathsf{Diag}_{a,b}\) (the PBW residual
of the Weyl/cyclic-trace comparison) lies entirely in this positive
part.  Under the radial-parts datum, its vanishing is proved by testing
on the Cartan
(Lemmas~\ref{lem:app-radial-cartan-moyal-representation}
and~\ref{lem:app-radial-dunkl-calogero-cancellation}) against the
Levasseur--Stafford radial-kernel datum.  A signed-row obstruction
\((\phi,-\lambda)\in\ker B_{a,b}^*\) (a covector pairing a
cyclic-word functional \(\phi\in\mathsf V_{a,b}^*\) with a
diagonal-target functional \(\lambda\in\mathsf{Diag}_{a,b}^*\)) of the
\(\lambda(E^+_{a,b})-\phi(T_{a,b})\neq0\) type is excluded under that
datum.

The formal Weyl/Moyal coefficients on the constant symplectic plane
and the Capelli triangular trace normalization on the matrix Weyl
algebra reduce the comparison to a finite-dimensional linear-algebra
problem in the free indexed normal-diagram algebra; finite-rank
Procesi--Razmyslov stable injectivity extends free identities to
all-\(N\) identities.  Harish-Chandra, Wallach, and
Levasseur--Stafford
\cite{harish-chandra-invariant-differential-operators}
\cite{wallach-reductive-invariant-differential}
\cite{levasseur-stafford-harish-chandra}
\cite{levasseur-stafford-kernel-harish-chandra} supply the source
literature for the invariant radial homomorphism and kernel theorem.
The exact radial input used here is not imported as an external
all-bidegree theorem; it is the three-clause supplied datum of
Statement~\ref{stmt:app-radial-parts-hypotheses}.  The conditional
all-bidegree criterion uses clauses~(1) and~(2) (Levasseur--Stafford
radial-kernel) and clause~(3) only for the specific Moyal-bracket
polynomial of bidegree \((a,b)\).  Pure-power clause~(3) is
tautological because Weyl-symmetric and normal-ordered traces of
pure powers coincide.  The reduced open-line Wick graph realization
with midpoint propagator \(\tfrac12\operatorname{sgn}(t-s)\)
realizes the same Moyal product to all orders and is the bridge to
the open-line BV sector.

The edge strips \(\mathfrak o_{2,b}=\mathfrak o_{a,2}=0\) close in
closed form (Proposition~\ref{prop:app-edge-pbw-telescoping}).  The
Procesi--Razmyslov reduction
(Proposition~\ref{prop:app-radial-pr-unconditional-window}) closes
the listed interior bidegrees by finite rational linear algebra; the
explicit list of bidegrees verified in this mode
consists of the non-edge bidegrees of total degree at most \(13\)
together with the six total-degree-\(14\) bidegrees
\[
  (3,11),(11,3),(4,10),(10,4),(5,9),(9,5),
\]
and is given in Table~\ref{tab:app-radial-finite-window-status}.
An unlisted bidegree remains governed by the obstruction class
\(\Omega^{\mathrm{rad}}_{a,b}\in\operatorname{coker}B_{a,b}\) unless
Statement~\ref{stmt:app-radial-parts-hypotheses} is supplied.  Reflection
symmetry \((a,b)\leftrightarrow(b,a)\) is implemented by the
symplectic involution \(z_1\leftrightarrow z_2\),
\(X\leftrightarrow Y\); both sides of the radial bicomplex are
equivariant under it.

\begin{table}[t]
\centering
\small
\begin{tabular}{@{}>{\raggedright\arraybackslash}p{0.30\linewidth}
  >{\raggedright\arraybackslash}p{0.19\linewidth}
  >{\raggedright\arraybackslash}p{0.43\linewidth}@{}}
\toprule
\textbf{Bidegree range} & \textbf{Status} & \textbf{Input and consequence} \\
\midrule
\(a=2\) or \(b=2\) &
proved &
closed-form edge telescoping; \(\Omega^{\mathrm{rad}}_{a,b}=0\) unconditionally \\
\(a,b\ge3,\ a+b\le13\) &
proved &
exact rational reduction over the indexed-diagram basis \\
\(a+b=14\), with
\(\begin{gathered}
(3,11),(11,3),(4,10),\\
(10,4),(5,9),(9,5)
\end{gathered}\) &
proved &
exact rational reduction over the indexed-diagram basis \\
\((6,8),(8,6)\), and the balanced case \((7,7)\) &
open finite systems &
classical cyclic boundary equations close; no decorated target solution
and no signed row in \(\ker B_{a,b}^*\) is certified \\
unlisted bidegree &
open unless supplied &
governed by
\(\Omega^{\mathrm{rad}}_{a,b}\in\operatorname{coker}B_{a,b}\)
unless the radial-parts datum or a per-bidegree rational reduction is supplied \\
\bottomrule
\end{tabular}
\caption{Verified and open finite windows for the radial bicomplex.}
\label{tab:app-radial-finite-window-status}
\end{table}

\begin{defn}[Formal Weyl/Moyal convention]
\label{def:app-radial-moyal-convention}
  Let
  \[
    P=\partial_{z_1}\otimes\partial_{z_2}
      -\partial_{z_2}\otimes\partial_{z_1}.
  \]
  The formal Weyl product on \(\C[[z_1,z_2]][[\hbar_{\mathrm{W}}]]\) is
  \[
    f*g=m\circ\exp\left(\frac{\hbar_{\mathrm{W}}}{2}P\right)(f\otimes g),
  \]
  and
  \[
    [f,g]_{\mathrm{raw}}=f*g-g*f,\qquad
    \{f,g\}_{\hbar_{\mathrm{W}}}=\hbar_{\mathrm{W}}^{-1}[f,g]_{\mathrm{raw}}.
  \]
  The reduced quantum Hamiltonian Lie algebra is the scalar quotient
  \[
    \bar A_{\hbar_{\mathrm{W}}}=\C[[z_1,z_2]][[\hbar_{\mathrm{W}}]]/\C[[\hbar_{\mathrm{W}}]]
  \]
  as a Lie algebra.  The construction uses this quotient only in the
  Lie-algebra sense.
\end{defn}

\begin{thm}[Exact formal Weyl/Moyal coefficients]
\label{thm:app-radial-moyal-coefficients}
  For \(f=z_1^kz_2^l\) and \(g=z_1^mz_2^n\),
  \[
    P^r(f,g)=
    \sum_{a=0}^r(-1)^a\binom ra
      (k)_{r-a}(l)_a(m)_a(n)_{r-a}
      z_1^{k+m-r}z_2^{l+n-r}.
  \]
  Consequently
  \[
    [f,g]_{\mathrm{raw}}
      =\sum_{\substack{r\geq1\\r\ \mathrm{odd}}}
        \frac{2}{2^r r!}\hbar_{\mathrm{W}}^rP^r(f,g),
  \]
  equivalently, for \(s\geq0\), the raw coefficient of
  \(\hbar_{\mathrm{W}}^{2s+1}\) is
  \[
    \frac{1}{2^{2s}(2s+1)!}\,P^{2s+1}(f,g),
  \]
  and the normalized Hamiltonian bracket has coefficient
  \[
    \frac{\hbar_{\mathrm{W}}^{2s}}{2^{2s}(2s+1)!}\,P^{2s+1}(f,g).
  \]
  The first and third raw commutator coefficients are
  \[
    \hbar_{\mathrm{W}} P(f,g),\qquad \frac{\hbar_{\mathrm{W}}^3}{24}P^3(f,g),
  \]
  and all even raw commutator coefficients vanish.  This is a theorem
  inside the formal Weyl algebra of the constant symplectic plane,
	independent of the finite-\(N\) radial-parts datum and Costello graph
  specialization.
\end{thm}

\begin{proof}
  The two summands in
  \(P=\partial_{z_1}\otimes\partial_{z_2}
     -\partial_{z_2}\otimes\partial_{z_1}\)
  commute as differential operators on the tensor product.  The
  binomial theorem gives
  \[
    P^r=\sum_{a=0}^r(-1)^a\binom ra
      \partial_{z_1}^{r-a}\partial_{z_2}^{a}
      \otimes
      \partial_{z_1}^{a}\partial_{z_2}^{r-a}.
  \]
  Applying this to the two monomials gives the falling-factorial
  formula.  Interchanging \(f\) and \(g\) changes \(P^r(f,g)\) by
  \((-1)^r\), so the commutator retains exactly the odd powers and
  contributes the coefficient \(2/(2^r r!)\) in order \(r\).
\end{proof}

\begin{cor}[All odd Moyal coefficients]
\label{cor:app-radial-all-odd-moyal-coefficients}
  The Weyl/Moyal commutator has only odd raw powers of \(\hbar_{\mathrm{W}}\).  Its
  odd raw and normalized coefficients are
  \[
    [f,g]_{\mathrm{raw}}^{(2s+1)}
      =
      \frac{\hbar_{\mathrm{W}}^{2s+1}}{2^{2s}(2s+1)!}P^{2s+1}(f,g),
    \qquad
    \{f,g\}_{\hbar_{\mathrm{W}}}^{(2s)}
      =
      \frac{\hbar_{\mathrm{W}}^{2s}}{2^{2s}(2s+1)!}P^{2s+1}(f,g).
  \]
  In particular the next two raw coefficients after the third-order
  term are
  \[
    \frac{\hbar_{\mathrm{W}}^5}{1920}P^5(f,g),\qquad
    \frac{\hbar_{\mathrm{W}}^7}{322560}P^7(f,g).
  \]
\end{cor}

\begin{proof}
  Put \(r=2s+1\) in
  Theorem~\ref{thm:app-radial-moyal-coefficients}.  Then
  \(2/(2^{2s+1}(2s+1)!)=1/(2^{2s}(2s+1)!)\).  Dividing by \(\hbar_{\mathrm{W}}\)
  gives the normalized coefficient.
\end{proof}

\begin{cor}[First and third formal graph weights]
\label{cor:app-formal-first-third-weights}
  In the ordered formal product
  \(m\circ\exp(\hbar_{\mathrm{W}} P/2)\), the one-cross-contraction term has weight
  \(\hbar_{\mathrm{W}} P(f,g)/2\), and the three-cross-contraction term has weight
  \(\hbar_{\mathrm{W}}^3P^3(f,g)/48\).  After antisymmetrising the ordered product,
  the commutator weights are
  \[
    \hbar_{\mathrm{W}} P(f,g),\qquad \frac{\hbar_{\mathrm{W}}^3}{24}P^3(f,g).
  \]
\end{cor}

\begin{proof}
  The ordered exponential contributes
  \[
    \frac{1}{r!}\left(\frac{\hbar_{\mathrm{W}}}{2}\right)^r P^r(f,g)
  \]
  in order \(r\).  For \(r=1\) this is \(\hbar_{\mathrm{W}} P(f,g)/2\); for \(r=3\)
  it is \(\hbar_{\mathrm{W}}^3P^3(f,g)/(3!\,2^3)=\hbar_{\mathrm{W}}^3P^3(f,g)/48\).  Since
  \(P(g,f)=-P(f,g)\) and \(P^3(g,f)=-P^3(f,g)\), subtracting the reverse
  ordered product doubles these two ordered contributions.
\end{proof}

\begin{ex}[A higher-order coefficient check]
\label{ex:app-radial-higher-moyal-check}
  For
  \[
    f=z_1^4z_2^2,\qquad g=z_1^2z_2^4,
  \]
  the first nonzero odd contractions are
  \[
    P(f,g)=12z_1^5z_2^5,\qquad
    P^3(f,g)=-960z_1^3z_2^3,\qquad
    P^5(f,g)=11520z_1z_2 .
  \]
  Hence
  \[
    [f,g]_{\mathrm{raw}}
      =
      12\hbar_{\mathrm{W}} z_1^5z_2^5
      -40\hbar_{\mathrm{W}}^3 z_1^3z_2^3
      +6\hbar_{\mathrm{W}}^5 z_1z_2 .
  \]
  The third coefficient is \(-960/24=-40\).  The fifth coefficient is
  \(11520/1920=6\).  This example checks the numerical normalization
  beyond the first nontrivial term, independently of radial parts.
\end{ex}

\begin{proof}
  Apply the falling-factorial formula of
  Theorem~\ref{thm:app-radial-moyal-coefficients}.  For \(r=1\),
  \[
    4\cdot4-2\cdot2=12.
  \]
  For \(r=3\),
  \[
  \begin{aligned}
    P^3(f,g)
    &=
      \bigl(24\cdot24
      -3\cdot12\cdot2\cdot2\cdot12
      +3\cdot4\cdot2\cdot2\cdot4\bigr)z_1^3z_2^3  \\
    &=-960z_1^3z_2^3 .
  \end{aligned}
  \]
  For \(r=5\), only the \(a=1\) and \(a=2\) summands survive, and they
  give
  \[
    \bigl(-5\cdot24\cdot2\cdot2\cdot24
      +10\cdot24\cdot2\cdot2\cdot24\bigr)z_1z_2
    =
    11520z_1z_2 .
  \]
  Substituting the odd Moyal coefficients
  \(1\), \(1/24\), and \(1/1920\) gives the displayed raw commutator.
\end{proof}

\begin{defn}[Matrix Weyl algebra and Weyl traces]
\label{def:app-radial-matrix-weyl}
  Let \(\mathcal W_N\) be the completed Weyl algebra generated by
  entries of two \(N\times N\) matrices \(X,Y\) with
  \[
    [X^i{}_j,Y^k{}_l]=\hbar_{\mathrm{W}}\,\delta^i_l\delta^k_j,\qquad
    [X^i{}_j,X^k{}_l]=[Y^i{}_j,Y^k{}_l]=0.
  \]
  The normal-ordered quantum comoment for conjugation is
  \[
    \widehat\mu_\epsilon=-\operatorname{Tr}([\epsilon,X]Y),
    \qquad
    \hbar_{\mathrm{W}}^{-1}[\widehat\mu_\epsilon,a]=\epsilon\cdot a.
  \]
  For \(f\in\C[z_1,z_2][[\hbar_{\mathrm{W}}]]\), define the Weyl-ordered trace
  \[
    J_N(f)=\operatorname{Tr}\operatorname{Op}_{\mathrm W}(f)(X,Y).
  \]
  The degree-zero quantum Hamiltonian reduction uses the effective
  traceless comoment ideal
  \[
    \mathrm N_{\mathcal W_N}
      (\mathcal W_N\widehat\mu(\lie{sl}_N))/
    \mathcal W_N\widehat\mu(\lie{sl}_N).
  \]
  The scalar Hamiltonian line survives this moment ideal; below it is
  removed by quotienting the empty trace before connected stable-trace
  extraction.
\end{defn}

\begin{rmk}[Weyl algebra relation versus finite-rank matrix equation]
\label{rmk:app-weyl-versus-matrix-equation}
The commutation relation in Definition~\ref{def:app-radial-matrix-weyl}
is an entrywise canonical commutation relation in the matrix Weyl
algebra.  It is not the finite-dimensional equation
\([X,Y]=\hbar_{\mathrm{W}}I_N\), and no pair of \(N\times N\) matrices
can satisfy that equation for \(\hbar_{\mathrm{W}}\ne0\), since the trace
of a commutator is zero.  Quantum Hamiltonian reduction by
\(\lie{sl}_N\) imposes the traceless normal-ordered comoment.  In the
normal ordering used here this comoment is represented by
\[
  M_N^{\mathrm{qmm}}
  =
  YX-XY+\hbar_{\mathrm{W}}N I_N,
  \qquad
  \operatorname{Tr}M_N^{\mathrm{qmm}}=0
\]
inside \(\mathcal W_N\).  Thus \(M_N^{\mathrm{qmm}}\equiv0\) modulo
\(\mathcal W_N\widehat\mu(\lie{sl}_N)\) means the traceless quantum
moment-map congruence.  The scalar \(\hbar_{\mathrm{W}}N\) is the
Capelli contact term surviving on trace generators before the empty
trace quotient, not a scalar value of a finite-rank commutator.
\end{rmk}

\begin{lem}[Capelli--Weyl triangular normalization]
\label{lem:app-radial-capelli-triangular}
  Let \(T_{a,b}=\operatorname{Tr}(X^aY^b)\) be the normal-ordered trace.
  In the quantum Hamiltonian reduction, the locally proved
  Weyl-normal-ordering formula is
  \[
    J_N(z_1^az_2^b)\equiv
      \sum_{r=0}^{\min(a,b)}
        \left(-\frac{\hbar_{\mathrm{W}} N}{2}\right)^r
        r!\binom ar\binom br\,T_{a-r,b-r}.
  \]
  The inverse triangular system is
  \[
    T_{a,b}\equiv
      \sum_{r=0}^{\min(a,b)}
        \left(\frac{\hbar_{\mathrm{W}} N}{2}\right)^r
        r!\binom ar\binom br\,J_N(z_1^{a-r}z_2^{b-r}).
  \]
\end{lem}

\begin{proof}
  Total Weyl symmetrisation is obtained from normal ordering by
  commuting \(Y\)'s past \(X\)'s.  Each contraction contributes the
  Capelli contact term coming from
  \[
    YX-XY+\hbar_{\mathrm{W}} N I_N\equiv0
      \pmod{\mathcal W_N\widehat\mu(\lie{sl}_N)}.
  \]
  Choosing \(r\) contracted \(X\)-slots and \(r\) contracted \(Y\)-slots
  gives \(\binom ar\binom br\), pairing them gives \(r!\), and Weyl
  ordering contributes \(2^{-r}\).  The sign is negative because the
  displayed traceless comoment congruence rewrites a normal-ordering swap as the
  lower triangular term \(-\hbar_{\mathrm{W}} N\).  The diagonal term is \(r=0\),
  hence the system is unipotent over \(\C[[\hbar_{\mathrm{W}} N]]\); inversion
  changes \(-\hbar_{\mathrm{W}} N/2\) to \(+\hbar_{\mathrm{W}} N/2\).

  The cited Capelli literature is the source for the
  Capelli contact phenomenon; the displayed contraction count fixes
  the bivariate trace normalization used here.
\end{proof}

\begin{stmt}[Radial-parts supplied datum]
\label{stmt:app-radial-parts-hypotheses}
  The finite-\(N\) radial-parts step separates two inputs.

  The first input is the invariant radial homomorphism and kernel
  theorem for the matrix differential operators, in the Harish-Chandra,
  Wallach, and Levasseur--Stafford source tradition cited above.  This
  supplied datum gives the
  Harish-Chandra map on invariants, the regular-Cartan landing, and the
  kernel statement for the quantum moment-map ideal.  Hence the quantum
  Hamiltonian reduction by conjugation injects into Weyl-invariant
  differential operators on the regular Cartan.  With
  \(\lambda_1,\ldots,\lambda_N\) the Cartan coordinates and
  \(\Delta=\prod_{i<j}(\lambda_i-\lambda_j)\), the radial-part map is
  conjugated by \(\Delta\).

  The second input is the supplied Weyl/Capelli trace-image datum of
  item~(3).  That datum is the matched normalization used in the
  formal-local comparison.
  The finite exact checks below test that normalization in ranks \(2\)
  and \(3\); the all-\(N\) Capelli/Weyl image formula remains a supplied
  datum.

  The supplied datum is used through the following three assertions.
  \begin{enumerate}
    \item The Vandermonde-conjugated radial map is a homomorphism from
      the invariant quantum Hamiltonian reduction by conjugation to
      Weyl-invariant differential operators on the regular Cartan.
    \item Its kernel is precisely the quantum moment-map ideal; hence it
      is injective on the reduced quotient.
    \item The Weyl trace generators used here have image
  \[
    \operatorname{Rad}_0(J_N(f))
      =\Delta^{-1}\left(\sum_{i=1}^N
        \operatorname{Op}_{\mathrm W}
        (f)(\lambda_i,-\hbar_{\mathrm{W}}\partial_{\lambda_i})\right)\Delta.
  \]
  \end{enumerate}
  The nonlinear-shear test in
  Lemma~\ref{lem:app-radial-nonlinear-shear-obstruction} shows that
  pure \(z_1\)- and \(z_2\)-power images together with classical
  symbol-level shear compatibility determine only the pure-power datum.
  Item~(3) requires a genuine quantum-corrected Egorov theorem in the
  Weyl algebra.
\end{stmt}

\begin{defn}[Radially normalized Weyl trace class]
Assume only the homomorphism and injectivity clauses of
Statement~\ref{stmt:app-radial-parts-hypotheses}.  For
\(f\in\C[z_1,z_2][[\hbar_{\mathrm{W}}]]\), a class \(J_N^{\mathrm{rad}}(f)\) in the
degree-zero quantum Hamiltonian reduction is called radially normalized
if
\[
  \operatorname{Rad}_0(J_N^{\mathrm{rad}}(f))
  =\Delta^{-1}S_N(f)\Delta .
\]
Existence is a normalization datum; uniqueness follows from radial
injectivity.  Statement~\ref{stmt:app-radial-parts-hypotheses}(3) is
exactly the assertion that the ordinary Weyl trace \(J_N(f)\) equals
\(J_N^{\mathrm{rad}}(f)\).
\end{defn}

\begin{defn}[Vandermonde gauge variables]
\label{def:app-radial-vandermonde-gauge-variables}
  Set
  \[
    \nabla_i^\Delta
      =
      \Delta^{-1}(-\hbar_{\mathrm{W}}\partial_{\lambda_i})\Delta
      =
      -\hbar_{\mathrm{W}}\partial_{\lambda_i}
      -\hbar_{\mathrm{W}}\sum_{j\neq i}\frac{1}{\lambda_i-\lambda_j}.
  \]
  The logarithmic summand is the type \(A\) Dunkl--Calogero correction
  in the Vandermonde gauge.  For \(f\in\C[z_1,z_2][[\hbar_{\mathrm{W}}]]\), put
  \[
    S_N(f)=\sum_{i=1}^N
      \operatorname{Op}_{\mathrm W}
      (f)(\lambda_i,-\hbar_{\mathrm{W}}\partial_{\lambda_i})
  \]
  and
  \[
    R_N(f)=
      \sum_{i=1}^N
      \operatorname{Op}_{\mathrm W}
      (f)(\lambda_i,\nabla_i^\Delta)
      =
      \Delta^{-1}S_N(f)\Delta .
  \]
\end{defn}

\begin{lem}[Dunkl--Calogero correction cancellation]
\label{lem:app-radial-dunkl-calogero-cancellation}
  The Vandermonde-gauge variables satisfy
  \[
    [\lambda_i,\lambda_j]=0,\qquad
    [\nabla_i^\Delta,\nabla_j^\Delta]=0,\qquad
    [\lambda_i,\nabla_j^\Delta]=\hbar_{\mathrm{W}}\,\delta_{ij}.
  \]
  Consequently
  \[
    \hbar_{\mathrm{W}}^{-1}[R_N(f),R_N(g)]=R_N(\{f,g\}_{\hbar_{\mathrm{W}}})
  \]
  for all \(f,g\).  In particular, all logarithmic
  Dunkl--Calogero terms introduced by expanding
  \(\Delta^{-1}S_N(-)\Delta\) cancel in the commutator identity.
\end{lem}

\begin{proof}
  The operator \(\nabla_i^\Delta\) is conjugate to
  \(-\hbar_{\mathrm{W}}\partial_{\lambda_i}\) by multiplication by \(\Delta\).
  The three displayed commutation relations are therefore the ordinary
  Weyl relations for
  \((\lambda_i,-\hbar_{\mathrm{W}}\partial_{\lambda_i})\), conjugated by
  \(\Delta\).  The correction terms are a pure gauge connection:
  their curvature is zero because the unconjugated partial derivatives
  commute.  Hence the \(i\)-th summand is again a one-particle Weyl
  representation of the formal Moyal algebra, and different summands
  commute.  This gives the displayed commutator identity.  The
  Vandermonde-conjugated form contains the Calogero potential.
\end{proof}

\begin{lem}[Radial-map injectivity on the reduced quotient]
\label{lem:app-radial-map-injectivity}
  Under Statement~\ref{stmt:app-radial-parts-hypotheses}, if an invariant
  Weyl operator \(D\) in the normalizer of the quantum moment-map ideal
  satisfies
  \[
    \operatorname{Rad}_0(D)=0,
  \]
  then \(D\) is zero in the degree-zero quantum Hamiltonian reduction.
\end{lem}

\begin{proof}
  Statement~\ref{stmt:app-radial-parts-hypotheses} identifies the kernel
  of the Vandermonde-conjugated Harish-Chandra radial map with the
  quantum moment-map ideal.  Passing to the normalizer quotient by that
  ideal gives an injective map into Weyl-invariant differential
  operators on the regular Cartan.  Thus an element with zero radial
  image represents the zero class in the reduced quotient.
\end{proof}

\begin{lem}[Nonlinear shear obstruction to weaker radial-parts data]
\label{lem:app-radial-nonlinear-shear-obstruction}
  Pure-power radial images plus classical symbol-level compatibility
  with nonlinear shears determine only the pure-power Weyl-trace image
  normalization.  In the one-particle Weyl algebra
  \([\lambda,y]=\hbar_{\mathrm{W}}\), the coefficient of \(t\) in
  \[
    (y+t\lambda^2)^4
  \]
  is
  \[
    4\lambda^2y^3-12\hbar_{\mathrm{W}}\lambda y^2+8\hbar_{\mathrm{W}}^2 y,
  \]
  whereas
  \[
    4\operatorname{Op}_{\mathrm W}(z_1^2z_2^3)
    =
    4\lambda^2y^3-12\hbar_{\mathrm{W}}\lambda y^2+6\hbar_{\mathrm{W}}^2 y.
  \]
  The difference \(2\hbar_{\mathrm{W}}^2y\) is non-scalar.  Thus
  Statement~\ref{stmt:app-radial-parts-hypotheses}(3) requires the quantum
  Egorov correction in addition to classical nonlinear shear covariance.
\end{lem}

\begin{proof}
  Expand to first order in \(t\):
  \[
    (y+t\lambda^2)^4
    =
    y^4+t\sum_{j=0}^3 y^j\lambda^2y^{3-j}+O(t^2).
  \]
  Using \(y\lambda=\lambda y-\hbar_{\mathrm{W}}\), the sum normal-orders to
  \(4\lambda^2y^3-12\hbar_{\mathrm{W}}\lambda y^2+8\hbar_{\mathrm{W}}^2y\).
  Weyl symmetrisation of \(z_1^2z_2^3\) gives
  \(\lambda^2y^3-3\hbar_{\mathrm{W}}\lambda y^2+\frac32\hbar_{\mathrm{W}}^2y\), hence the
  displayed value after multiplication by \(4\).  The discrepancy is a
  nonconstant one-particle operator, hence survives the scalar
  Hamiltonian quotient.
\end{proof}

\begin{prop}[Quantum shear congruence suffices for the all-\(N\) image]\label{prop:app-quantum-shear-suffices-radial-image}
  Assume the homomorphism and injectivity clauses of
  Statement~\ref{stmt:app-radial-parts-hypotheses}, and assume the pure
  power images
  \[
    \operatorname{Rad}_0(J_N(z_1^r))=\Delta^{-1}S_N(z_1^r)\Delta,
    \qquad
    \operatorname{Rad}_0(J_N(z_2^r))=\Delta^{-1}S_N(z_2^r)\Delta
  \]
  for all \(r\geq0\).  Suppose moreover that, for all \(a,b\geq0\), the
  quantum shear congruence
  \[
    \hbar_{\mathrm{W}}^{-1}\bigl[
      J_N(z_1^{a+1}/(a+1)),J_N(z_2^{b+1})
    \bigr]
    \equiv
    J_N\!\left(
      \{z_1^{a+1}/(a+1),z_2^{b+1}\}_{\hbar_{\mathrm{W}}}
    \right)
    \pmod {I_N}
  \]
  holds in the degree-zero quantum Hamiltonian reduction, where \(I_N\)
  is the quantum moment-map ideal.  Then the exact Weyl-trace image in
  Statement~\ref{stmt:app-radial-parts-hypotheses}(3) holds for every
  monomial \(z_1^az_2^b\), hence for every polynomial \(f\).
\end{prop}

\begin{proof}
  The normalized Moyal bracket gives
  \[
    \{z_1^{a+1}/(a+1),z_2^{b+1}\}_{\hbar_{\mathrm{W}}}
    =
    (b+1)z_1^az_2^b
    +\sum_{s\ge1}
      \frac{\hbar_{\mathrm{W}}^{2s}(a+1)_{2s+1}(b+1)_{2s+1}}
      {(a+1)2^{2s}(2s+1)!}
      z_1^{a-2s}z_2^{b-2s}.
  \]
  Every summand after the leading term has lower total degree.  The pure
  power cases start the induction.  Assume the asserted radial image is
  known in all total degrees \(<a+b\).  Applying the radial homomorphism
  to the quantum shear congruence identifies the radial image of the
  left-hand commutator with the Moyal commutator of the already known
  pure-power radial images.  Lemma~\ref{lem:app-radial-cartan-moyal-representation}
  computes this as the Vandermonde-conjugated Weyl quantization of the
  displayed Moyal bracket.  By the induction hypothesis all lower-degree
  terms have the required images, so the leading coefficient
  \((b+1)\) determines the image of \(J_N(z_1^az_2^b)\).  Radial
  injectivity then identifies the class in the reduced quotient.
\end{proof}

\begin{lem}[Moment-ideal primitive form of the quantum shear datum]
\label{lem:app-quantum-shear-primitive-obstruction}
  In the notation of
  Proposition~\ref{prop:app-quantum-shear-suffices-radial-image}, the
  quantum shear congruence is equivalent to the following explicit
  left-ideal membership problem.  Put
  \[
    E_{a,b,N}
    =
    \sum_{j=0}^b\operatorname{Tr}(Y^jX^aY^{b-j})
    -
    J_N\!\left(
      \{z_1^{a+1}/(a+1),z_2^{b+1}\}_{\hbar_{\mathrm{W}}}
    \right).
  \]
  Then the congruence holds if and only if
  \(E_{a,b,N}\in\mathcal W_N\widehat\mu(\lie{sl}_N)\).  Equivalently,
  after choosing normal-ordered matrix generators
  \(\widehat\mu^i{}_j\) for the traceless quantum moment-map ideal, one
  must construct coefficients \(A^j{}_i(a,b,N)\in\mathcal W_N\) such
  that
  \[
    E_{a,b,N}=\sum_{i,j}A^j{}_i(a,b,N)\widehat\mu^i{}_j .
  \]
\end{lem}

\begin{proof}
  Since the pure powers are already Weyl ordered,
  \(J_N(z_1^{a+1}/(a+1))=\operatorname{Tr}(X^{a+1})/(a+1)\) and
  \(J_N(z_2^{b+1})=\operatorname{Tr}(Y^{b+1})\).  The Weyl relations give
  \[
    \hbar_{\mathrm{W}}^{-1}
    \bigl[\operatorname{Tr}(X^{a+1})/(a+1),Y\bigr]
    =X^a .
  \]
  Applying the Leibniz rule to \(Y^{b+1}\) gives
  \[
    \hbar_{\mathrm{W}}^{-1}
    \bigl[
      J_N(z_1^{a+1}/(a+1)),J_N(z_2^{b+1})
    \bigr]
    =
    \sum_{j=0}^b\operatorname{Tr}(Y^jX^aY^{b-j}).
  \]
  Subtracting the \(J_N\)-image of the normalized Moyal bracket gives
  exactly \(E_{a,b,N}\).  Reduction modulo \(I_N\) is reduction modulo
  the left ideal
  \(\mathcal W_N\widehat\mu(\lie{sl}_N)\), so the displayed congruence
  is precisely the asserted membership.  The coefficient formula is the
  same membership written after a choice of generators for that left
  ideal.
\end{proof}

\begin{stmt}[Free trace-diagram quantum-lift obstruction]
\label{stmt:app-quantum-shear-trace-diagram-obstruction}
  Put \(M=YX-XY+\hbar_{\mathrm{W}} N I\), and interpret
  \[
    \operatorname{Tr}(W M)=\sum_{i,j}W^j{}_iM^i{}_j
  \]
  as a left-ideal expression.  Let \(\mathsf C_{a,b}\) be the rational
  vector space spanned by words \(W\) with \(a-1\) letters \(X\) and
  \(b-1\) letters \(Y\), and let \(\mathsf{Cyc}_{a,b}\) be the cyclic
  word space with \(a\) letters \(X\) and \(b\) letters \(Y\).  The
  leading term of such a primitive is the cyclic boundary equation
  \[
    \partial c=
    \sum_{j=0}^b[Y^jX^aY^{b-j}]
    -\frac{b+1}{\binom{a+b}{a}}
      \sum_{U\in \operatorname{Sh}(a,b)}[U],
    \qquad
    \partial W=[WYX]-[WXY].
  \]
  An all-\(N\) primitive is stronger.  It requires equality after free
  indexed normal ordering, including every contraction term which splits
  a trace into multi-traces:
  \[
    \mathcal N_{\mathrm{free}}\!\left(
      \sum_W c_W\operatorname{Tr}(W M)\right)
    =
    \mathcal N_{\mathrm{free}}(E_{a,b,N}).
  \]
  The cyclic equation is the associated-graded part of this quantum
  equation.

  In bidegree \((a,b)=(4,4)\), the lexicographic pivot solution of the
  cyclic equation is
  \[
  \begin{aligned}
    C^{\mathrm{cl}}_{4,4}
    &=
      \frac{31}{7}\operatorname{Tr}(XXXYYY\,M)
      -\frac{4}{7}\operatorname{Tr}(XXYXYY\,M)
      -\frac{4}{7}\operatorname{Tr}(XXYYXY\,M)\\
    &\quad
      +\frac{27}{7}\operatorname{Tr}(XXYYYX\,M)
      +\frac{23}{7}\operatorname{Tr}(XYXXYY\,M)
      +\frac{1}{7}\operatorname{Tr}(XYXYXY\,M)\\
    &\quad
      +\operatorname{Tr}(XYXYYX\,M)
      -\frac{2}{7}\operatorname{Tr}(XYYXXY\,M)
      +\frac{15}{7}\operatorname{Tr}(XYYXYX\,M).
  \end{aligned}
  \]
  Its quantum residual is nonzero.  In rank \(N=2\), exact normal
  ordering gives
  \[
  \begin{aligned}
    R^{\mathrm{cl}}_{4,4,2}
    &:=
      E_{4,4,2}-C^{\mathrm{cl}}_{4,4}\\
    &=
      \frac{43}{70}\hbar_{\mathrm{W}}^2\bigl(
        4\operatorname{Tr}(XXYY)-5\operatorname{Tr}(XYXY)
        +2\operatorname{Tr}(XYYX)\\
    &\hspace{4.6cm}
        -5\operatorname{Tr}(YXYX)
        +4\operatorname{Tr}(YYXX)
      \bigr)\neq0 .
  \end{aligned}
  \]
  After expanding the displayed traces into normal-ordered matrix
  entries, the coefficient of
  \(\hbar_{\mathrm{W}}^2X^1{}_1X^1{}_1Y^1{}_2Y^2{}_1\) is \(43/7\).

  The residual is killed by the classical-kernel correction
  \[
  \begin{aligned}
    K_{4,4}
    &=
    \frac{43}{14}\bigl(
      \operatorname{Tr}(XXYXYY\,M)
      -\operatorname{Tr}(XYXXYY\,M)
      \\
    &\qquad
      -\operatorname{Tr}(XYYXYX\,M)
      +\operatorname{Tr}(XYYYXX\,M)
    \bigr),
  \end{aligned}
  \]
  whose cyclic boundary is zero.  The free normal-form computation below
  proves that this correction cancels the \((4,4)\) residual before any
  finite-rank trace identity is imposed.  The decorated PBW Stokes
  identity closes uniformly: by
  Theorem~\ref{thm:app-radial-bicomplex-acyclic}, for every bidegree
  \[
    \Omega^{\mathrm{rad}}_{a,b}=0
    \quad\Longleftrightarrow\quad
    D^\square_{a,b}H_{a,b}=R^{\mathrm{free}}_{a,b},
    \qquad
    D^\square_{a,b}=C^+_{a,b}\partial_2 ,
  \]
  under Statement~\ref{stmt:app-radial-parts-hypotheses}.
  Unconditionally, the
  bidegrees in Remark~\ref{rmk:app-radial-finite-verification} are
  closed by Procesi--Razmyslov verification
  (Proposition~\ref{prop:app-radial-pr-unconditional-window}); the
  normalized signed-dual-row obstruction set is empty there.
\end{stmt}

\begin{proof}
  The leading \(\hbar_{\mathrm{W}}^0\) component of
  \(\operatorname{Tr}(W M)\) is
  \([WYX]-[WXY]\) in the cyclic trace quotient; comparing it with the
  leading term of \(E_{a,b,N}\) gives the displayed cyclic equation.
  Normal ordering in the indexed Weyl algebra contracts every occurrence
  of a \(Y\)-entry moved past an \(X\)-entry.  These contractions identify
  matrix indices and may split one cyclic trace into a product of cyclic
  traces.  Hence equality in the cyclic quotient is only the associated
  graded condition.

  The \((4,4)\) coefficients displayed above solve the cyclic equation
  in the stated pivot gauge.  Direct exact normal ordering in
  \(\mathcal W_2\) gives the displayed residual.  Its nonvanishing follows
  from the indicated normal monomial coefficient.  Applying \(\partial\)
  to the four words in \(K_{4,4}\) gives two copies of each resulting
  cyclic word with opposite signs, so \(K_{4,4}\in\ker\partial\).  The
  finite computation proves the obstruction to the uncorrected cyclic
  lift.  The all-\(N\) identity for \(K_{4,4}\) is given by the free
  trace-diagram normal form below.
\end{proof}

\begin{defn}[Free indexed normal diagrams]
\label{def:app-free-indexed-normal-diagrams}
  Let \(u=L_1\cdots L_m\) be a word in \(X,Y\).  Write the trace
  \(\operatorname{Tr}(u)\) as the directed \(m\)-gon with vertices
  \(v_0,\ldots,v_{m-1}\) and edge \(e_r:v_r\to v_{r+1}\), labelled
  \(L_r\), with indices read modulo \(m\).  A contraction matching
  \(\Gamma\) is a partial matching of pairs
  \[
    p<q,\qquad L_p=Y,\qquad L_q=X .
  \]
  For each pair \((p,q)\in\Gamma\), delete the two edges \(e_p,e_q\)
  and identify
  \[
    v_q\sim v_{p+1},\qquad v_p\sim v_{q+1}.
  \]
  Let \(e(u,\Gamma)\) be the number of isolated vertices left after
  these deletions and identifications.  Let
  \(\mathsf D(u,\Gamma)\) be the remaining directed labelled graph,
  with all \(X\)-edges placed before all \(Y\)-edges in the ambient
  commutative normal-ordered algebra of matrix entries.  Disconnected
  components are retained; they are the split-trace terms.  The empty
  graph is denoted \(1\).
\end{defn}

\begin{thm}[Free indexed normal-ordering formula]
\label{thm:app-free-indexed-normal-ordering}
  In the universal indexed matrix Weyl algebra with
  \[
    Y^k{}_lX^i{}_j=X^i{}_jY^k{}_l-\hbar_{\mathrm{W}}\,\delta^i_l\delta^k_j,
  \]
  the normal-ordered trace word is
  \[
    \mathcal N_{\mathrm{free}}\operatorname{Tr}(u)
    =
    \sum_{\Gamma}
      (-\hbar_{\mathrm{W}})^{|\Gamma|}
      N^{e(u,\Gamma)}
      \mathsf D(u,\Gamma),
  \]
  where the sum runs over all contraction matchings
  \(\Gamma\) of \(Y\)-before-\(X\) inversions in \(u\).  Consequently
  \[
  \begin{aligned}
    \mathcal N_{\mathrm{free}}\operatorname{Tr}(W M)
    &=
      \mathcal N_{\mathrm{free}}\operatorname{Tr}(WYX)
      -\mathcal N_{\mathrm{free}}\operatorname{Tr}(WXY)\\
    &\quad
      +\hbar_{\mathrm{W}} N\,\mathcal N_{\mathrm{free}}\operatorname{Tr}(W),
  \end{aligned}
  \]
  with \(M=YX-XY+\hbar_{\mathrm{W}} N I\), interpreted as the left-ideal expression
  \(\operatorname{Tr}(W M)=\sum_{i,j}W^j{}_iM^i{}_j\).
\end{thm}

\begin{proof}
  Move the \(X\)-entries leftward through the \(Y\)-entries by the PBW
  relation displayed above.  Choosing the commutator summand at a swap
  selects one pair \(Y_p,X_q\), contributes the factor \(-\hbar_{\mathrm{W}}\), and
  imposes exactly the two Kronecker identifications
  \(v_q\sim v_{p+1}\) and \(v_p\sim v_{q+1}\).  Choosing the ordered
  summand leaves both edges and continues the PBW ordering.  Since a
  generator can be contracted only once, the possible choices are
  precisely the partial matchings of original inversions.  After all
  \(X\)-entries have been moved left, the uncontracted entries form the
  normal diagram \(\mathsf D(u,\Gamma)\).  Each isolated index loop
  contributes a free summation over \(1,\ldots,N\), hence \(N\).  This
  proves the formula for trace words.  The formula for
  \(\operatorname{Tr}(W M)\) is the defining expansion of \(M\), with the
  scalar matrix term contributing \(\hbar_{\mathrm{W}} N\operatorname{Tr}(W)\).
\end{proof}

\begin{lem}[Stable injectivity for normal trace diagrams]
\label{lem:app-free-trace-diagram-stable-injectivity}
  Fix \(d\).  Let \(\mathsf{Diag}_{\le d}\) be the rational span of free
  indexed normal diagrams with at most \(d\) labelled nonempty edges,
  with coefficients in \(\Q[\hbar_{\mathrm{W}},N]\).  The joint specialization map
  to \(N\times N\) matrices for all ranks \(N\ge d\) is injective.  Thus
  a free normal-form equality in \(\mathsf{Diag}_{\le d}\) is an all-\(N\)
  matrix Weyl identity.  A nonzero free residual remains visible until
  the stable Procesi--Razmyslov range is reached.
\end{lem}

\begin{proof}
  Polarize in the labelled edges.  A linear relation among polarized
  closed directed trace diagrams of edge degree at most \(d\) is a trace
  identity for generic matrices of degree at most \(d\).  The
  Procesi--Razmyslov second fundamental theorem says that the trace
  identities are generated by the fundamental antisymmetrizer on
  \(N+1\) index lines.  Hence every identity of degree at most
  \(d\) is zero in ranks \(N\ge d\).  Since the coefficients are
  polynomials in the formal scalar \(N\), vanishing in every rank
  \(N\ge d\) implies that each coefficient polynomial vanishes.  Depolarizing
  gives the stated injectivity.
\end{proof}

\begin{prop}[First free residual and its kernel correction]
\label{prop:app-free-trace-diagram-K44}
  Use the normal-diagram notation of
  Definition~\ref{def:app-free-indexed-normal-diagrams}; for instance
  \(\mathsf D(X_{0\to1}Y_{1\to0})\) denotes the directed two-edge
  normal diagram with one \(X\)-edge and one \(Y\)-edge.  Then
  \[
  \begin{aligned}
    \mathcal N_{\mathrm{free}}
      (E_{4,4,N}-C^{\mathrm{cl}}_{4,4})
    &=
      \frac{43}{7}\hbar_{\mathrm{W}}^2\Bigl(
      \mathsf D(X_{0\to1}X_{1\to2}Y_{2\to3}Y_{3\to0})\\
    &\quad
      -\mathsf D(X_{0\to1}X_{2\to3}Y_{1\to2}Y_{3\to0})\\
    &\quad
      -\hbar_{\mathrm{W}}\,\mathsf D(X_{0\to0}\mid Y_{0\to0})
      +\hbar_{\mathrm{W}} N\,\mathsf D(X_{0\to1}Y_{1\to0})
      \Bigr).
  \end{aligned}
  \]
  Moreover
  \[
    \mathcal N_{\mathrm{free}}(K_{4,4})
    =
    \mathcal N_{\mathrm{free}}
      (E_{4,4,N}-C^{\mathrm{cl}}_{4,4}),
  \]
  and therefore
  \[
    \mathcal N_{\mathrm{free}}
      (E_{4,4,N}-C^{\mathrm{cl}}_{4,4}-K_{4,4})=0.
  \]
  This is an all-\(N\) identity in the free indexed normal-diagram
  algebra, proved before any rank specialization.
\end{prop}

\begin{proof}
  Apply Theorem~\ref{thm:app-free-indexed-normal-ordering} to the nine
  trace words in \(C^{\mathrm{cl}}_{4,4}\) and to the Weyl-trace
  expression defining \(E_{4,4,N}\).  The finite matching enumeration is
  rational and takes place in the free diagram algebra.  It leaves exactly
  the four displayed normal diagrams.  Applying the same formula to the
  four trace words in \(K_{4,4}\) gives the same four-diagram vector.
  Exact enumeration in the free trace-diagram basis gives the
  four-term residual above and gives zero residual after adding
  \(K_{4,4}\).
\end{proof}

\begin{defn}[Kernel-correction residual and decorated necklace complex]
\label{def:app-free-kernel-correction-obstruction}
  Let \(\mathsf W_{a,b}\) be the rational span of words with \(a-1\)
  letters \(X\) and \(b-1\) letters \(Y\), and let
  \(\mathsf V_{a,b}=\mathsf{Cyc}_{a,b}\) be the rational span of cyclic
  words with \(a\) letters \(X\) and \(b\) letters \(Y\).  Define the
  two-colour necklace multigraph \(G_{a,b}\) by taking
  \(\mathsf V_{a,b}\) as its vertex span and by assigning to
  \(W\in\mathsf W_{a,b}\) the oriented edge
  \[
    e_W:[WXY]\longrightarrow[WYX].
  \]
  Its incidence map is
  \[
    \partial\colon\mathsf W_{a,b}\to\mathsf V_{a,b},\qquad
    \partial W=[WYX]-[WXY],
  \]
  and \(Z_{a,b}:=\ker\partial\) is the cycle space.

  Let \(\mathsf{Diag}_{a,b}\) denote the positive free indexed
  normal-diagram target in bidegree \((a,b)\): the uncontracted
  \(\hbar_{\mathrm{W}}^0\) cyclic component is removed and all PBW contraction,
  split-trace, and Capelli \(\hbar_{\mathrm{W}} N\)-terms are retained.  Write
  \(\pi_+\) for this projection.  Define the decorated contraction
  cochain by
  \[
    C^+_{a,b}(W)
    =
    \pi_+\mathcal N_{\mathrm{free}}\operatorname{Tr}(WM),
    \qquad
    M=YX-XY+\hbar_{\mathrm{W}} NI.
  \]
  The correction map is
  \[
    \mathcal C_{a,b}=A_{a,b}:=
    C^+_{a,b}|_{Z_{a,b}}\colon
    Z_{a,b}\longrightarrow\mathsf{Diag}_{a,b}.
  \]

  Put
  \[
    T_{a,b}=
    \sum_{j=0}^b[Y^jX^aY^{b-j}]
    -\frac{b+1}{\binom{a+b}{a}}
      \sum_{U\in\operatorname{Sh}(a,b)}[U].
  \]
  For any solution \(c^{\mathrm{cl}}_{a,b}\) of
  \(\partial c=T_{a,b}\), define
  \[
    E^+_{a,b}
    =
    \pi_+\mathcal N_{\mathrm{free}}(E_{a,b,N})
  \]
  and put
  \[
    R_{a,b}^{\mathrm{free}}
    =
    E^+_{a,b}-C^+_{a,b}(c^{\mathrm{cl}}_{a,b})
  \]
  and define
  \[
    \mathfrak o_{a,b}
    =
    [R_{a,b}^{\mathrm{free}}]\in
    \operatorname{coker}\mathcal C_{a,b}.
  \]
  This is the kernel-correction class after a choice of classical lift.
  The radial obstruction itself is the lift-independent class
  \(\Omega^{\mathrm{rad}}_{a,b}\) and the signed-row criterion of
  Proposition~\ref{prop:app-radial-dual-potential-obstruction}.
\end{defn}

\begin{stmt}[Decorated kernel-correction Hodge criterion]
\label{stmt:app-free-kernel-homotopy-obstruction}
  Give \(Z_{a,b}\) and \(\mathsf{Diag}_{a,b}\) the inner products for
  which the chosen word-cycle and free normal-diagram bases are
  orthonormal.  Let \(A=A_{a,b}\) and define
  \[
    \Delta^{\mathrm{dec}}_{a,b}=AA^*,\qquad
    \mathsf H^{\mathrm{dec}}_{a,b}=\ker A^*
      =(\operatorname{im}A)^\perp .
  \]
  Let
  \[
    \Pi^{\mathrm{harm}}_{a,b}\colon
    \mathsf{Diag}_{a,b}\to\mathsf H^{\mathrm{dec}}_{a,b}
  \]
  be the orthogonal projection.  For every \(a,b\geq1\), the following
  conditions are equivalent.
  \begin{enumerate}
    \item There exists \(K_{a,b}\in Z_{a,b}\) such that
      \[
        C^+_{a,b}(K_{a,b})=R^{\mathrm{free}}_{a,b}.
      \]
    \item The kernel-correction class \(\mathfrak o_{a,b}\) vanishes.
    \item The decorated harmonic projection vanishes:
      \[
        \Pi^{\mathrm{harm}}_{a,b}R^{\mathrm{free}}_{a,b}=0.
      \]
    \item Every functional
      \(\lambda\in\mathsf{Diag}_{a,b}^{*}\) with
      \[
        \lambda C^+_{a,b}(K)=0\qquad(K\in Z_{a,b})
      \]
      also satisfies \(\lambda(R^{\mathrm{free}}_{a,b})=0\).
  \end{enumerate}
  When these equivalent conditions hold, the canonical Hodge correction is
  \[
    K_{a,b}=
    A_{a,b}^{*}(\Delta^{\mathrm{dec}}_{a,b})^{-1}
    R^{\mathrm{free}}_{a,b},
  \]
  with the inverse taken on \(\operatorname{im}A_{a,b}\).  Thus the
  fixed-lift kernel-correction equation is equivalent to
  \[
    \Pi^{\mathrm{harm}}_{a,b}R^{\mathrm{free}}_{a,b}=0
    \qquad\text{in the chosen bidegree.}
  \]
  The lift-independent obstruction
  \(\Omega^{\mathrm{rad}}_{a,b}\) and its normalized signed-row form
  are given by
  Proposition~\ref{prop:app-radial-dual-potential-obstruction}.
\end{stmt}

\begin{proof}
  If \(c\) and \(c'\) solve the cyclic equation, then
  \(c-c'\in Z_{a,b}\).  Hence
  \(R(c')-R(c)=-A_{a,b}(c'-c)\), so the class in
  \(\operatorname{coker}A_{a,b}\), and its harmonic projection, are
  independent of the chosen cyclic lift.

  The finite-dimensional Hodge decomposition for \(A\) gives
  \[
    \mathsf{Diag}_{a,b}=\operatorname{im}A\oplus\ker A^* .
  \]
  Therefore \(R^{\mathrm{free}}_{a,b}\in\operatorname{im}A\) if and
  only if its projection to \(\ker A^*\) is zero.  This proves the
  equivalence of the correction equation, the intermediate cokernel
  class, and the harmonic projection.  The functional criterion is the
  annihilator form of the same statement:
  \(R\in\operatorname{im}A\) if and only if every
  \(\lambda\in(\operatorname{im}A)^\perp=\ker A^*\) vanishes on \(R\).
  If \(R\in\operatorname{im}A\), then
  \(AA^*(AA^*)^{-1}R=R\) on \(\operatorname{im}A\), so the displayed
  \(K_{a,b}\) is a correction.  Conversely, any correction places
  \(R^{\mathrm{free}}_{a,b}\) in \(\operatorname{im}A\).
\end{proof}

\begin{defn}[Square-cell PBW Stokes complex]
\label{def:app-square-cell-pbw-stokes-complex}
  Let \(\mathsf S_{a,b}\) be the rational span of square cells obtained
  from two commuting adjacent transpositions in a cyclic word with \(a\) letters
  \(X\) and \(b\) letters \(Y\).  A based representative has the form
  \[
    uXYvXYw,
  \]
  and its oriented square is
  \[
    uXYvXYw\to uYXvXYw\to uYXvYXw\to
    uXYvYXw\to uXYvXYw .
  \]
  Each side is read as the corresponding edge
  \(e_W:[WXY]\to[WYX]\) of the necklace graph.  This defines
  \[
    \partial_2\colon \mathsf S_{a,b}\longrightarrow\mathsf W_{a,b}.
  \]
  The decorated PBW Stokes map is
  \[
    D^\square_{a,b}
      =
      C^+_{a,b}\partial_2\colon
      \mathsf S_{a,b}\longrightarrow\mathsf{Diag}_{a,b}.
  \]
\end{defn}

\begin{lem}[Ordinary necklace Stokes theorem]
\label{lem:app-ordinary-necklace-stokes}
  For \(a,b\geq1\),
  \[
    \operatorname{im}\partial_2=\ker\partial .
  \]
\end{lem}

\begin{proof}
  Work first with based words.  A binary word with \(a\) letters \(X\)
  and \(b\) letters \(Y\) is the same as an order ideal in the rectangle
  poset \([a]\times[b]\), via its inversion set.  An adjacent
  \(XY\leftrightarrow YX\) move toggles one extremal box.  The cubical
  complex of the distributive lattice of order ideals, with squares
  generated by commuting toggles, is contractible.  Hence its rational
  one-cycles are generated by square boundaries.  Passing to cyclic words
  is the quotient by a finite cellular rotation action; over \(\Q\),
  transfer identifies the quotient rational homology with the invariant
  rational homology upstairs.  Thus every necklace cycle is a rational
  sum of square boundaries.
\end{proof}

\begin{prop}[Dual-potential form of the radial obstruction]
\label{prop:app-radial-dual-potential-obstruction}
  Define the finite-dimensional map
  \[
    B_{a,b}\colon\mathsf W_{a,b}\longrightarrow
    \mathsf V_{a,b}\oplus\mathsf{Diag}_{a,b},
    \qquad
    B_{a,b}(W)=(\partial W,C^+_{a,b}(W)),
  \]
  and the target
  \[
    \eta_{a,b}=(T_{a,b},E^+_{a,b}).
  \]
  Let
  \[
    \Omega^{\mathrm{rad}}_{a,b}
    =
    [\eta_{a,b}]\in\operatorname{coker}B_{a,b}.
  \]
  For fixed \((a,b)\), the following are equivalent.
  \begin{enumerate}
    \item \(\Omega^{\mathrm{rad}}_{a,b}=0\).
    \item For one, equivalently every, classical lift
      \(c^{\mathrm{cl}}_{a,b}\) with
      \(\partial c^{\mathrm{cl}}_{a,b}=T_{a,b}\), there is an
      \(H_{a,b}\in\mathsf S_{a,b}\) such that
      \[
        D^\square_{a,b}H_{a,b}=R^{\mathrm{free}}_{a,b}.
      \]
    \item For one, equivalently every, classical lift
      \(c^{\mathrm{cl}}_{a,b}\), the residual
      \(R^{\mathrm{free}}_{a,b}\) lies in
      \(\operatorname{im}(C^+_{a,b}|_{\ker\partial})\).
    \item \(\Pi^{\mathrm{harm}}_{a,b}R^{\mathrm{free}}_{a,b}=0\).
    \item For every pair
      \[
        (\phi,\lambda)\in
        \mathsf V_{a,b}^*\oplus\mathsf{Diag}_{a,b}^*
      \]
      satisfying
      \[
        \lambda C^+_{a,b}(W)=\phi(\partial W)
        \qquad(W\in\mathsf W_{a,b}),
      \]
      one has
      \[
        \lambda(E^+_{a,b})=\phi(T_{a,b}).
      \]
  \end{enumerate}
  Thus
  \[
    \Omega^{\mathrm{rad}}_{a,b}=0
    \quad\Longleftrightarrow\quad
    D^\square_{a,b}H_{a,b}=R^{\mathrm{free}}_{a,b}
  \]
  for one, equivalently every, classical lift.  The obstruction in
  bidegree \((a,b)\) is exactly a signed dual row
  \[
    (\phi,-\lambda)\in\ker B_{a,b}^*
  \]
  with
  \[
    \lambda(E^+_{a,b})-\phi(T_{a,b})\ne0.
  \]
  Equivalently, after rescaling, it is a normalized row
  \[
    \lambda\in\ker(D^\square_{a,b})^*,
    \qquad
    \lambda(R^{\mathrm{free}}_{a,b})=1 .
  \]
\end{prop}

\begin{proof}
  The pair \((\phi,\lambda)\) in the statement is the signed dual row
  \((\phi,-\lambda)\) on
  \(\mathsf V_{a,b}\oplus\mathsf{Diag}_{a,b}\).  It annihilates
  \(B_{a,b}(W)\) precisely when
  \(\lambda C^+_{a,b}(W)=\phi(\partial W)\).  Since the spaces are finite
  dimensional,
  \[
    \eta_{a,b}\in\operatorname{im}B_{a,b}
    \quad\Longleftrightarrow\quad
    \ell(\eta_{a,b})=0
    \text{ for every }\ell\in\ker B_{a,b}^* .
  \]
  Evaluating \((\phi,-\lambda)\) on \(\eta_{a,b}\) gives
  \(\phi(T_{a,b})-\lambda(E^+_{a,b})\), proving the equivalence of
  \(\Omega^{\mathrm{rad}}_{a,b}=0\) with the signed-row condition.

  If \(c^{\mathrm{cl}}_{a,b}\) is a classical lift, then
  \[
    \eta_{a,b}-B_{a,b}(c^{\mathrm{cl}}_{a,b})
    =
    (0,R^{\mathrm{free}}_{a,b}).
  \]
  Adding \(B_{a,b}(K)\) with \(K\in\ker\partial\) changes only the
  decorated diagram component by \(C^+_{a,b}(K)\).  Hence
  \(\eta_{a,b}\in\operatorname{im}B_{a,b}\) if and only if
  \(R^{\mathrm{free}}_{a,b}\in
  \operatorname{im}(C^+_{a,b}|_{\ker\partial})\).  By
  Lemma~\ref{lem:app-ordinary-necklace-stokes},
  \(\ker\partial=\operatorname{im}\partial_2\), so this image is exactly
  \(\operatorname{im}D^\square_{a,b}\).  The Hodge equivalence is the
  decomposition of
  Statement~\ref{stmt:app-free-kernel-homotopy-obstruction}.

  Finally, finite-dimensional duality for \(D^\square_{a,b}\) identifies
  the cokernel of \(D^\square_{a,b}\) with \(\ker(D^\square_{a,b})^*\).
  A nonzero residual class therefore has a row
  \(\lambda\in\ker(D^\square_{a,b})^*\) with
  \(\lambda(R^{\mathrm{free}}_{a,b})\ne0\).  Rescale the row to make the
  value \(1\).  The equality
  \(\lambda C^+_{a,b}(W)=\phi(\partial W)\) identifies this row with the
  signed row \((\phi,-\lambda)\in\ker B_{a,b}^*\), and evaluating that
  signed row on \(\eta_{a,b}\) gives
  \(\phi(T_{a,b})-\lambda(E^+_{a,b})\).
\end{proof}

\begin{rmk}[Potential form and the bare graph Green operator]
\label{rmk:app-decorated-hodge-potential-form}
  A functional \(\lambda\in\mathsf{Diag}_{a,b}^{*}\) satisfies
  \(\lambda C^+_{a,b}|_{\ker\partial}=0\) precisely when the edge cochain
  \[
    W\longmapsto \lambda C^+_{a,b}(W)
  \]
  has zero periods on the necklace graph \(G_{a,b}\).  The graph is
  connected for \(a,b>0\), since adjacent \(XY\leftrightarrow YX\)
  moves connect all two-colour necklaces.  Hence this zero-period
  condition is equivalent to a vertex potential
  \(\phi_\lambda\in\mathsf V_{a,b}^{*}\), unique up to a constant, with
  \[
    \lambda C^+_{a,b}(W)
    =
    \phi_\lambda([WYX])-\phi_\lambda([WXY])
    =
    \phi_\lambda(\partial W).
  \]
  Because \(T_{a,b}\) has total coefficient zero, \(\phi_\lambda(T_{a,b})\)
  is independent of the additive constant.  The radial obstruction
  vanishes, equivalently \(\Omega^{\mathrm{rad}}_{a,b}=0\), if and only
  if every such \(\lambda\) satisfies the potential identity
  \[
    \lambda\pi_+\mathcal N_{\mathrm{free}}(E_{a,b,N})
    =
    \phi_\lambda(T_{a,b}).
  \]
  By Lemma~\ref{lem:app-ordinary-necklace-stokes}, the same row condition
  is \(\lambda D^\square_{a,b}=0\), and
  \[
    \lambda(R^{\mathrm{free}}_{a,b})
    =
    \lambda E^+_{a,b}
    -\lambda C^+_{a,b}(c^{\mathrm{cl}}_{a,b})
    =
    \lambda E^+_{a,b}
    -\phi_\lambda(T_{a,b}).
  \]

  The ordinary graph Green operator for the incidence map
  \(\partial\) is therefore insufficient.  It constructs a classical
  cyclic lift \(c^{\mathrm{cl}}_{a,b}\) of \(T_{a,b}\), i.e. it solves
  only the uncontracted \(\hbar_{\mathrm{W}}^0\) incidence equation.  The quantum
  correction equation is instead
  \(A_{a,b}K_{a,b}=R^{\mathrm{free}}_{a,b}\), equivalently
  \(D^\square_{a,b}H_{a,b}=R^{\mathrm{free}}_{a,b}\), inside the positive
  PBW diagram target, where contraction diagrams, split traces, and the
  Capelli term are visible.  A functorial right inverse from arbitrary
  positive diagram residuals to necklace edge cochains is a separate
  datum \(S^{\square}_{a,b}\) with
  \(D^\square_{a,b}S^{\square}_{a,b}=\operatorname{id}\) on the chosen
  residual subcomplex.
  The finite-dimensional Green operator used here is the decorated
  square-cell one attached to
  \(D^\square_{a,b}(D^\square_{a,b})^{*}\).
\end{rmk}

\begin{defn}[Radial bicomplex and PBW \(\hbar_{\mathrm{W}}\)-filtration]
\label{def:app-radial-bicomplex}
  Fix \(a,b\geq1\).  The \emph{radial bicomplex} in bidegree \((a,b)\)
  is the augmented diagram
  \[
    \mathcal R_{a,b}\colon\quad
    \mathsf S_{a,b}
      \xrightarrow{\partial_2}
    \mathsf W_{a,b}
      \xrightarrow{\partial}
    \mathsf V_{a,b},
    \qquad
    \mathsf W_{a,b}\xrightarrow{C^+_{a,b}}\mathsf{Diag}_{a,b}
  \]
  with \(\partial,\partial_2,C^+_{a,b}\) as in
  Definitions~\ref{def:app-free-kernel-correction-obstruction}
  and~\ref{def:app-square-cell-pbw-stokes-complex}.  Equip
  \(\mathsf{Diag}_{a,b}\) with the \emph{contraction filtration}
  \[
    F^p\mathsf{Diag}_{a,b}=\bigoplus_{p'\geq p}\mathsf{Diag}_{a,b}^{(p')},
  \]
  where \(\mathsf{Diag}_{a,b}^{(p)}\) is the rational span of free
  indexed normal diagrams with exactly \(p\) PBW contractions.  The
  Capelli scalar contributes one unit of \(\hbar_{\mathrm{W}}\) and is recorded in
  \(\mathsf{Diag}_{a,b}^{(0)}\) with an explicit factor \(\hbar_{\mathrm{W}} N\).
  Consequently \(C^+_{a,b}\) and its restriction
  \(D^\square_{a,b}=C^+_{a,b}\partial_2\) take values in
  \(F^1\mathsf{Diag}_{a,b}\), and the residual
  \(R^{\mathrm{free}}_{a,b}\) lies in \(F^1\mathsf{Diag}_{a,b}\).  The
  graded pieces \(\mathrm{gr}^p\mathsf{Diag}_{a,b}\) carry the Cartan
  one-particle Moyal action of
  Lemma~\ref{lem:app-radial-cartan-moyal-representation}.
\end{defn}

\begin{thm}[Radial bicomplex acyclicity]
\label{thm:app-radial-bicomplex-acyclic}
  Assume the radial-parts datum of
  Statement~\ref{stmt:app-radial-parts-hypotheses} in full.
  Then for every interior bidegree \((a,b)\) with \(a,b\geq1\),
  \[
    \Omega^{\mathrm{rad}}_{a,b}=0,
    \qquad
    D^\square_{a,b}H_{a,b}=R^{\mathrm{free}}_{a,b}
    \text{ has a solution }H_{a,b}\in\mathsf S_{a,b}.
  \]
  Equivalently the radial bicomplex \(\mathcal R_{a,b}\) is acyclic
  at the residual class \(R^{\mathrm{free}}_{a,b}\) on the decorated
  side: every element of \(\mathsf{Diag}_{a,b}\) of the form
  \(R^{\mathrm{free}}_{a,b}\) lies in
  \(\operatorname{im}D^\square_{a,b}\).  Equivalently, the normalized
  signed-dual-row obstruction set is empty.
\end{thm}

\begin{proof}
  The argument proceeds in three steps.  Step~1 uses only clause~(1)
  (radial homomorphism) and the Cartan-one-particle Moyal
  homomorphism of
  Lemma~\ref{lem:app-radial-cartan-moyal-representation}; Step~2
  uses the mixed-monomial image clause~(3) of
  Statement~\ref{stmt:app-radial-parts-hypotheses} for the
  Moyal-bracket polynomial of bidegree \((a,b)\), which is the
  Weyl/Capelli normalization datum; Step~3 uses clause~(2) (Levasseur--Stafford
  radial-kernel theorem) and stable Procesi--Razmyslov injectivity.

  \emph{Step 1 (Cartan-side commutator).}  Compute the
  Vandermonde-conjugated Harish-Chandra radial image of the matrix
  Weyl commutator
  \(\hbar_{\mathrm{W}}^{-1}[J_N(z_1^{a+1}/(a+1)),J_N(z_2^{b+1})]\).  Pure-power
  Weyl traces equal normal-ordered traces, so
  \(J_N(z_1^{a+1})=\operatorname{Tr}(X^{a+1})\) and
  \(J_N(z_2^{b+1})=\operatorname{Tr}(Y^{b+1})\) with zero contractions;
  clause~(3) for pure powers is therefore tautological:
  \[
    \operatorname{Rad}_0(J_N(z_1^{a+1}))
      =\Delta^{-1}S_N(z_1^{a+1})\Delta
      =\sum_{i=1}^N\lambda_i^{a+1},
    \qquad
    \operatorname{Rad}_0(J_N(z_2^{b+1}))
      =\Delta^{-1}S_N(z_2^{b+1})\Delta
      =\sum_{i=1}^N(\nabla_i^\Delta)^{b+1}.
  \]
  Both pre-images are \(GL_N\)-invariant, so by clause~(1) the
  radial map is a homomorphism on their commutator.  Different Cartan
  summands commute by the third commutation relation of
  Lemma~\ref{lem:app-radial-dunkl-calogero-cancellation}, hence
  \[
    \operatorname{Rad}_0\bigl(
      \hbar_{\mathrm{W}}^{-1}[J_N(z_1^{a+1}/(a+1)),J_N(z_2^{b+1})]
    \bigr)
    =
    \sum_{i=1}^N
      \hbar_{\mathrm{W}}^{-1}\bigl[\lambda_i^{a+1}/(a+1),(\nabla_i^\Delta)^{b+1}\bigr].
  \]
  Lemma~\ref{lem:app-radial-cartan-moyal-representation} identifies
  each summand with the Cartan one-particle Weyl quantization of the
  normalized Moyal bracket:
  \[
    \hbar_{\mathrm{W}}^{-1}\bigl[\lambda_i^{a+1}/(a+1),(\nabla_i^\Delta)^{b+1}\bigr]
    =
    \operatorname{Op}_{\mathrm W}\!\bigl(
      \{z_1^{a+1}/(a+1),z_2^{b+1}\}_{\hbar_{\mathrm{W}}}
    \bigr)(\lambda_i,\nabla_i^\Delta).
  \]
  Summing over \(i\) and using the displayed formula in
  Definition~\ref{def:app-radial-vandermonde-gauge-variables} gives
  \[
    \operatorname{Rad}_0\bigl(
      \hbar_{\mathrm{W}}^{-1}[J_N(z_1^{a+1}/(a+1)),J_N(z_2^{b+1})]
    \bigr)
    =
    \Delta^{-1}S_N\!\bigl(
      \{z_1^{a+1}/(a+1),z_2^{b+1}\}_{\hbar_{\mathrm{W}}}
    \bigr)\Delta .
  \]
  Up to this point only the pure-power part of clause~(3) has entered.

  \emph{Step 2 (Spherical-reduction equality).}  By
  Statement~\ref{stmt:app-radial-parts-hypotheses}~clause~(3)
  applied to the Moyal-bracket polynomial
  \(\{z_1^{a+1}/(a+1),z_2^{b+1}\}_{\hbar_{\mathrm{W}}}\),
  \[
    \operatorname{Rad}_0\bigl(
      J_N(\{z_1^{a+1}/(a+1),z_2^{b+1}\}_{\hbar_{\mathrm{W}}})
    \bigr)
    =
    \Delta^{-1}S_N\!\bigl(
      \{z_1^{a+1}/(a+1),z_2^{b+1}\}_{\hbar_{\mathrm{W}}}
    \bigr)\Delta .
  \]
  Subtracting the two displayed identities,
  \(\operatorname{Rad}_0(E_{a,b,N})=0\).

  \emph{Step 3 (Universal lift via Procesi--Razmyslov).}  By
  Statement~\ref{stmt:app-radial-parts-hypotheses}~clause~(2)
  (Levasseur--Stafford radial-kernel theorem), the kernel of the
  Vandermonde-conjugated Harish-Chandra map on
  \(GL_N\)-invariant operators is exactly the moment-map ideal
  \(\mathcal W_N\widehat\mu(\lie{sl}_N)\).  Since
  \(E_{a,b,N}\) is \(GL_N\)-invariant and has zero radial image,
  \[
    E_{a,b,N}\in\mathcal W_N\widehat\mu(\lie{sl}_N)\qquad
    (\text{all }N\geq1).
  \]
  By
  Lemma~\ref{lem:app-quantum-shear-primitive-obstruction}, the
  membership \(E_{a,b,N}\in\mathcal W_N\widehat\mu(\lie{sl}_N)\)
  produces coefficients \(A^j{}_i(N)\in\mathcal W_N\) with
  \(E_{a,b,N}=\sum_{i,j}A^j{}_i(N)\widehat\mu^i{}_j\); the
  coefficients depend on \(N\) and on \(\hbar_{\mathrm{W}}\) a priori.  Apply the
  \(GL_N\)-symmetrization average to the lift; \(E_{a,b,N}\) is
  \(GL_N\)-invariant, so the symmetrization preserves the equality and
  produces \(GL_N\)-equivariant tensor coefficients.  By the Procesi
  classical fundamental theorem on trace ring generators (which
  embeds in stable Procesi--Razmyslov,
  Lemma~\ref{lem:app-free-trace-diagram-stable-injectivity}), every
  \(GL_N\)-equivariant tensor coefficient with values in
  \(\mathcal W_N\) of total edge degree \(\leq a+b\) is a
  \(\mathbb Q[\hbar_{\mathrm{W}}]\)-linear combination of word matrix-elements
  \(W^j{}_i\).  Polarizing in \(N\) reduces the resulting trace-form
  identity
  \[
    E_{a,b,N}=\sum_{W\in\mathsf W_{a,b}}c_W(a,b)\operatorname{Tr}(WM)
  \]
  to a finite-dimensional rational linear system in
  \(c=(c_W)_{W\in\mathsf W_{a,b}}\in\mathbb Q^{|\mathsf W_{a,b}|}\),
  decomposed into independent homogeneous components by the
  contraction filtration of
  Definition~\ref{def:app-radial-bicomplex} and the Capelli factor.
  The system is consistent in every homogeneous component because,
  for every \(N\geq a+b\), the matrix-Weyl-rank specialization has a
	  solution; stable Procesi--Razmyslov injectivity in those degrees
	  extends the rank-\(N\) consistency to free-side rational
  consistency.  Choose any \(\mathbb Q\)-rational solution
  \(c=\sum_W c_W W\in\mathsf W_{a,b}\).

  Project the resulting free identity through the cyclic-word target.
  At contraction degree \(0\),
  \(\sum_W c_W \pi_0\mathcal N_{\mathrm{free}}\operatorname{Tr}(WM)
   =\sum_W c_W([WYX]-[WXY])
   =\partial(\sum_W c_W W)\), so the underlying word
  \(c\in\mathsf W_{a,b}\) is a classical lift,
  \(\partial c=T_{a,b}\).  Subtracting any preferred classical lift
  \(c^{\mathrm{cl}}_{a,b}\), the difference
  \(K_{a,b}=c-c^{\mathrm{cl}}_{a,b}\) is a cycle in
  \(\ker\partial\), and on the decorated side
  \(\pi_+\mathcal N_{\mathrm{free}}(\sum_W c_W\operatorname{Tr}(WM))
   =E^+_{a,b}=C^+_{a,b}(c^{\mathrm{cl}}_{a,b})+C^+_{a,b}(K_{a,b})\),
  so \(C^+_{a,b}(K_{a,b})=R^{\mathrm{free}}_{a,b}\).  By the
  ordinary necklace Stokes
  Lemma~\ref{lem:app-ordinary-necklace-stokes}, every cycle in
  \(\ker\partial\) lies in
  \(\operatorname{im}\partial_2\), so there exists
  \(H_{a,b}\in\mathsf S_{a,b}\) with
  \(\partial_2H_{a,b}=K_{a,b}\); hence
  \(D^\square_{a,b}H_{a,b}
   =C^+_{a,b}\partial_2H_{a,b}
   =C^+_{a,b}(K_{a,b})
   =R^{\mathrm{free}}_{a,b}\).
  By Proposition~\ref{prop:app-radial-dual-potential-obstruction},
  \(\Omega^{\mathrm{rad}}_{a,b}=0\) and the normalized
  signed-dual-row obstruction set is empty.
\end{proof}

\begin{rmk}[Cartan one-particle Moyal lift]
\label{rmk:app-radial-bicomplex-mechanism}
  The proof of Theorem~\ref{thm:app-radial-bicomplex-acyclic} runs
  the same Cartan one-particle reduction in each bidegree.  At
  contraction degree~\(0\), the residual is the cyclic shear sum
  minus the Weyl average; the necklace Stokes
  Lemma~\ref{lem:app-ordinary-necklace-stokes} is the assertion that
  it is a coboundary in \(\partial\).  At each higher contraction
  degree~\(p\geq1\), the residual is the \(p\)-contraction component
  of the matrix-Weyl shear identity \emph{after} subtracting the
  classical lift; the Vandermonde-conjugated radial map evaluates
  this on the Cartan as a sum of one-particle Moyal commutators
  \(\hbar_{\mathrm{W}}^{-1}[\lambda_i^{a+1}/(a+1),(\nabla_i^\Delta)^{b+1}]\), and
  Lemma~\ref{lem:app-radial-cartan-moyal-representation} replaces
  each by its Weyl quantization of the Moyal-bracket polynomial.
  Statement~\ref{stmt:app-radial-parts-hypotheses}~clause~(3) for
  that polynomial closes the equality with
  \(\operatorname{Rad}_0(J_N(\{...\}_{\hbar_{\mathrm{W}}}))\); clause~(2) places the
  matrix-Weyl difference in the moment ideal.  Stable
  Procesi--Razmyslov polarization extracts the universal cycle
  \(K_{a,b}\), and the necklace Stokes lemma gives the square-cell
  primitive.  The contraction-filtered dimension of
  \(\mathrm{gr}^{\geq1}\mathsf{Diag}_{a,b}\) grows polynomially in
  \((a,b)\), matching the empirical
  \(K_{a,b}\)-correction support count in
  Remark~\ref{rmk:app-radial-finite-verification}.
\end{rmk}

\begin{cor}[Bidegree obstruction equals the clause-(3) defect on the Moyal-bracket polynomial]
\label{cor:app-radial-bicomplex-clause3-equivalence}
  Assume Statement~\ref{stmt:app-radial-parts-hypotheses}~clauses~(1)
  and~(2) (Levasseur--Stafford radial homomorphism and radial-kernel
  theorem).  For each bidegree \((a,b)\) with \(a,b\geq1\), the
  following are equivalent.
  \begin{enumerate}
    \item \(\Omega^{\mathrm{rad}}_{a,b}=0\).
    \item Statement~\ref{stmt:app-radial-parts-hypotheses}~clause~(3)
      holds modulo the moment ideal for the Moyal-bracket polynomial
      \(\{z_1^{a+1}/(a+1),z_2^{b+1}\}_{\hbar_{\mathrm{W}}}\):
      \[
        \operatorname{Rad}_0\!\bigl(
          J_N(\{z_1^{a+1}/(a+1),z_2^{b+1}\}_{\hbar_{\mathrm{W}}})
        \bigr)
        \equiv
        \Delta^{-1}S_N\!\bigl(
          \{z_1^{a+1}/(a+1),z_2^{b+1}\}_{\hbar_{\mathrm{W}}}
        \bigr)\Delta
        \pmod{I_N}.
      \]
  \end{enumerate}
  Equivalently, the bidegree-local cokernel obstruction is exactly
  the spherical-reduction defect of the radial image on the
  Moyal-bracket polynomial.
\end{cor}

\begin{proof}
  Step 1 of the proof of
  Theorem~\ref{thm:app-radial-bicomplex-acyclic} shows that, under
  clauses~(1) and~(2),
  \[
    \operatorname{Rad}_0\!\bigl(
      \hbar_{\mathrm{W}}^{-1}[J_N(z_1^{a+1}/(a+1)),J_N(z_2^{b+1})]
    \bigr)
    =
    \Delta^{-1}S_N\!\bigl(
      \{z_1^{a+1}/(a+1),z_2^{b+1}\}_{\hbar_{\mathrm{W}}}
    \bigr)\Delta .
  \]
  Subtracting the radial image of \(J_N(\{...\}_{\hbar_{\mathrm{W}}})\),
  \[
    \operatorname{Rad}_0(E_{a,b,N})
    =
    \Delta^{-1}S_N\!\bigl(
      \{...\}_{\hbar_{\mathrm{W}}}
    \bigr)\Delta
    -
    \operatorname{Rad}_0\!\bigl(J_N(\{...\}_{\hbar_{\mathrm{W}}})\bigr) .
  \]
  By clause~(2), \(E_{a,b,N}\in I_N\) iff
  \(\operatorname{Rad}_0(E_{a,b,N})=0\), iff the displayed clause-(3)
  identity holds.  Lemma~\ref{lem:app-quantum-shear-primitive-obstruction}
  identifies \(E_{a,b,N}\in I_N\) with
  \(\Omega^{\mathrm{rad}}_{a,b}=0\) via
  Proposition~\ref{prop:app-radial-dual-potential-obstruction}; the
  intermediate Procesi--Razmyslov reduction is in Step 3 of the
  bicomplex acyclicity proof.
\end{proof}

\begin{prop}[Procesi--Razmyslov exact verification window]
\label{prop:app-radial-pr-unconditional-window}
  For each bidegree \((a,b)\) and each \(N\geq a+b\), the
  finite-rank moment-ideal membership
  \(E_{a,b,N}\in\mathcal W_N\widehat\mu(\lie{sl}_N)\) holds if and
  only if the free indexed normal-diagram identity
  \(\mathcal N_{\mathrm{free}}(E_{a,b,N}-K_{a,b})=0\) holds for some
  cycle \(K_{a,b}\in\ker\partial\), and the latter is independent of
  \(N\).  Hence the all-\(N\) closure for fixed \((a,b)\) reduces to a
  finite rational linear-algebra problem on the
  \((\mathsf W_{a,b}\oplus\mathsf S_{a,b})\)-side, decided by exact
  rational normal-form calculation.  This verification
  closes \(\Omega^{\mathrm{rad}}_{a,b}=0\) by exact rational
  normal-form computation, independently of
  Statement~\ref{stmt:app-radial-parts-hypotheses}, for the
  bidegrees in
  Remark~\ref{rmk:app-radial-finite-verification}.
\end{prop}

\begin{proof}
  Stable Procesi--Razmyslov injectivity
  (Lemma~\ref{lem:app-free-trace-diagram-stable-injectivity}) makes
  the free indexed normal-diagram algebra a faithful joint model
  of all matrix Weyl algebras of rank \(N\geq a+b\).  Hence a free
  identity holds if and only if it holds in \(\mathcal W_N\) for every
  \(N\geq a+b\).  The kernel-correction equation
  \(\partial K_{a,b}=0\),
  \(\mathcal N_{\mathrm{free}}(E_{a,b,N}-C^+_{a,b}(K_{a,b}))
    =\mathcal N_{\mathrm{free}}(C^+_{a,b}(c^{\mathrm{cl}}_{a,b}))\)
  is finite-dimensional rational linear algebra on
  \((\mathsf W_{a,b}\oplus\mathsf S_{a,b})\) and gives the coefficient
  array \(K_{a,b}\), then the Stokes primitive
  \(H_{a,b}=\partial_2^{-1}K_{a,b}\) on the necklace square-cell
  complex by Lemma~\ref{lem:app-ordinary-necklace-stokes}.  Existence
  of \(K_{a,b}\) is therefore an exact finite computation in rational
  arithmetic.
\end{proof}

\begin{prop}[Edge PBW telescoping primitive]
\label{prop:app-edge-pbw-telescoping}
  For every \(b\geq2\), set
  \[
    C^{\mathrm{edge}}_{2,b}
    =
    \sum_{r=0}^{\lfloor(b-2)/2\rfloor}
      (b-1-2r)\operatorname{Tr}(Y^rXY^{b-1-r}M).
  \]
  For every \(a\geq2\), set
  \[
    C^{\mathrm{edge}}_{a,2}
    =
    \frac3{a+1}
    \sum_{r=0}^{\lfloor(a-2)/2\rfloor}
      (a-1-2r)\operatorname{Tr}(X^{a-1-r}YX^rM).
  \]
  Then the free indexed normal-diagram identities
  \[
    \mathcal N_{\mathrm{free}}
      (E_{2,b,N}-C^{\mathrm{edge}}_{2,b})=0,
    \qquad
    \mathcal N_{\mathrm{free}}
      (E_{a,2,N}-C^{\mathrm{edge}}_{a,2})=0
  \]
  hold for all displayed edge bidegrees.  Hence
  \[
    \mathfrak o_{2,b}=0,\qquad
    \mathfrak o_{a,2}=0
  \]
  with zero separate kernel correction in this edge gauge.
\end{prop}

\begin{proof}
  Take the \((2,b)\) strip.  Put
  \[
    \Theta_s^{(b)}=\mathsf D(XY^sXY^{b-s}),
    \qquad 0\leq s\leq b,
  \]
  and
  \[
    \Phi_s^{(b)}
    =
    \mathsf D(XY^s)\mid\mathsf D(Y^{b-1-s}),
    \qquad 0\leq s\leq b-1,
  \]
  where the vertical bar denotes disjoint union and
  \(\mathsf D(Y^0)\) is the isolated index loop, hence contributes the
  scalar factor \(N\).  For
  \[
    W_r=Y^rXY^{b-1-r},\qquad 0\leq r\leq b-1,
  \]
  Theorem~\ref{thm:app-free-indexed-normal-ordering} gives the PBW
  column identity
  \[
    \mathcal N_{\mathrm{free}}\operatorname{Tr}(W_rM)
    =
    \Theta_r^{(b)}-\Theta_{r+1}^{(b)}
    -\hbar_{\mathrm{W}}\Phi_r^{(b)}
    +\hbar_{\mathrm{W}}\Phi_{b-1-r}^{(b)}.
  \]
  Expand
  \[
    \operatorname{Tr}(W_rM)
    =
    \operatorname{Tr}(W_rYX)
    -\operatorname{Tr}(W_rXY)
    +\hbar_{\mathrm{W}} N\operatorname{Tr}(W_r).
  \]
  PBW matchings in the first two trace words away from the
  adjacent crossing given by \(M\) cancel in pairs.  The matchings
  using that adjacent crossing are cancelled, including their secondary
  contractions, by the scalar term \(\hbar_{\mathrm{W}} N\operatorname{Tr}(W_r)\).
  The uncontracted part leaves the boundary
  \(\Theta_r^{(b)}-\Theta_{r+1}^{(b)}\).  The two endpoint contraction
  families are the left and reflected right blocks, giving
  \(-\hbar_{\mathrm{W}}\Phi_r^{(b)}\) and \(+\hbar_{\mathrm{W}}\Phi_{b-1-r}^{(b)}\).  These are
  exactly the surviving terms.

  The target \(E_{2,b,N}\) contains only the first and third Moyal
  terms:
  \[
    J_N\!\left(
      \{z_1^3/3,z_2^{b+1}\}_{\hbar_{\mathrm{W}}}
    \right)
    =
    (b+1)J_N(z_1^2z_2^b)
    +\frac{(b+1)b(b-1)}{12}\hbar_{\mathrm{W}}^2\operatorname{Tr}(Y^{b-2}).
  \]
  Counting PBW matchings in the shear sum and in the Weyl average gives
  \[
    \mathcal N_{\mathrm{free}}(E_{2,b,N})
    =
    \sum_{r=0}^{\lfloor(b-2)/2\rfloor}
      (b-1-2r)
      \left(
        \Theta_r^{(b)}-\Theta_{r+1}^{(b)}
        -\hbar_{\mathrm{W}}\Phi_r^{(b)}
        +\hbar_{\mathrm{W}}\Phi_{b-1-r}^{(b)}
      \right).
  \]
  In contraction degree \(0\), this is the discrete boundary coefficient
  of the two-\(X\) cyclic separation.  In contraction degree \(1\), the
  coefficient of \(\Phi_s^{(b)}\) is
  \(-(b+1)+2(s+1)=2s-(b-1)\), the antisymmetric half-sum above.  In
  contraction degree \(2\), the difference between the shear count and
  the first-Moyal Weyl count is exactly
  \[
    \frac{(b+1)b(b-1)}{12}\,
    \mathcal N_{\mathrm{free}}\operatorname{Tr}(Y^{b-2}),
  \]
  which is cancelled by the third Moyal term in the displayed target.
  Combining this target expansion with the column identity proves
  \(\mathcal N_{\mathrm{free}}(E_{2,b,N}-C^{\mathrm{edge}}_{2,b})=0\).

  The \((a,2)\) strip is the same one-moving-letter calculation with
  \(Y\) moving through an \(X\)-string.  For
  \(V_r=X^{a-1-r}YX^r\),
  \[
    \mathcal N_{\mathrm{free}}\operatorname{Tr}(V_rM)
    =
    \widetilde\Theta_r^{(a)}-\widetilde\Theta_{r+1}^{(a)}
    -\hbar_{\mathrm{W}}\widetilde\Phi_r^{(a)}
    +\hbar_{\mathrm{W}}\widetilde\Phi_{a-1-r}^{(a)}.
  \]
  Here
  \[
    \widetilde\Theta_s^{(a)}
      =\mathsf D(X^{a-s}YX^sY),
    \qquad
    \widetilde\Phi_s^{(a)}
      =\mathsf D(X^{a-1-s})\mid\mathsf D(X^sY),
  \]
  with \(\mathsf D(X^0)\) again interpreted as the isolated index loop.
  The first and third Moyal coefficients for
  \(\{z_1^{a+1}/(a+1),z_2^3\}_{\hbar_{\mathrm{W}}}\) are
  \[
    3,\qquad \frac{a(a-1)}4,
  \]
  namely \(3/(a+1)\) times the corresponding \((2,a)\) coefficients.
  The same count therefore gives the stated
  \(C^{\mathrm{edge}}_{a,2}\) and proves the second identity.
\end{proof}

\begin{prop}[Free residual vanishings beyond \((4,4)\)]
\label{prop:app-free-residual-vanishings}
  In the free indexed normal-diagram algebra,
  \(\mathfrak o_{3,5}=\mathfrak o_{5,3}=0\), and the lexicographic
  cyclic lifts already have zero quantum residual.  Explicitly,
  \[
  \begin{aligned}
    E_{3,5,N}
    &=
      \frac{36}{7}\operatorname{Tr}(XXYYYY\,M)
      +\frac{24}{7}\operatorname{Tr}(XYXYYY\,M)\\
    &\quad
      -\frac{6}{7}\operatorname{Tr}(XYYXYY\,M)
      -\frac{6}{7}\operatorname{Tr}(XYYYXY\,M)\\
    &\quad
      +\frac{30}{7}\operatorname{Tr}(XYYYYX\,M)
      +\frac{12}{7}\operatorname{Tr}(YXYXYY\,M),
  \end{aligned}
  \]
  and
  \[
  \begin{aligned}
    E_{5,3,N}
    &=
      \frac{24}{7}\operatorname{Tr}(XXXXYY\,M)
      -\frac{4}{7}\operatorname{Tr}(XXXYXY\,M)\\
    &\quad
      +\frac{20}{7}\operatorname{Tr}(XXXYYX\,M)
      +\frac{12}{7}\operatorname{Tr}(XXYXXY\,M)\\
    &\quad
      -\frac{8}{7}\operatorname{Tr}(XXYXYX\,M)
      +\frac{16}{7}\operatorname{Tr}(XXYYXX\,M).
  \end{aligned}
  \]
  In total degree \(9\) the same free computation gives
  \(\mathfrak o_{3,6}=\mathfrak o_{6,3}=0\).  The cyclic lifts already
  have zero free residual:
  \[
  \begin{aligned}
    E_{3,6,N}
    &=
      \frac{25}{4}\operatorname{Tr}(XXYYYYY\,M)
      +\frac{19}{4}\operatorname{Tr}(XYXYYYY\,M)\\
    &\quad
      -\frac{3}{4}\operatorname{Tr}(XYYXYYY\,M)
      -\frac{3}{4}\operatorname{Tr}(XYYYXYY\,M)\\
    &\quad
      -\frac{3}{4}\operatorname{Tr}(XYYYYXY\,M)
      +\frac{11}{2}\operatorname{Tr}(XYYYYYX\,M)\\
    &\quad
      +\frac{13}{4}\operatorname{Tr}(YXYXYYY\,M)
      +\frac{1}{4}\operatorname{Tr}(YXYYXYY\,M)\\
    &\quad
      +\frac{7}{4}\operatorname{Tr}(YXYYYXY\,M),
  \end{aligned}
  \]
  and
  \[
  \begin{aligned}
    E_{6,3,N}
    &=
      \frac{25}{7}\operatorname{Tr}(XXXXXYY\,M)
      -\frac{3}{7}\operatorname{Tr}(XXXXYXY\,M)\\
    &\quad
      +\frac{22}{7}\operatorname{Tr}(XXXXYYX\,M)
      -\frac{3}{7}\operatorname{Tr}(XXXYXXY\,M)\\
    &\quad
      -\frac{6}{7}\operatorname{Tr}(XXXYXYX\,M)
      +\frac{19}{7}\operatorname{Tr}(XXXYYXX\,M)\\
    &\quad
      +\frac{16}{7}\operatorname{Tr}(XXYXXXY\,M)
      +\frac{1}{7}\operatorname{Tr}(XXYXXYX\,M)\\
    &\quad
      +\operatorname{Tr}(XXYXYXX\,M).
  \end{aligned}
  \]
\end{prop}

\begin{proof}
  Apply Theorem~\ref{thm:app-free-indexed-normal-ordering}.  For
  \((3,5)\) and \((5,3)\), the displayed coefficients solve the cyclic
  equation and the free normal-form residual is zero.  The
  total-degree-\(9\) formulas for \((3,6)\) and \((6,3)\) are checked in
  the same way; their residual is also zero, so \(K_{3,6}=K_{6,3}=0\)
  in this pivot gauge.  Exact normal-form enumeration gives the
  \((3,5)\), \((5,3)\), \((3,6)\), \((6,3)\), \((2,7)\), and \((7,2)\)
  cases and separately checks the edge theorem.
  The calculation takes place in the free diagram basis before
  finite-rank trace identities enter.
\end{proof}

\begin{lem}[Cartan one-particle Moyal representation]
\label{lem:app-radial-cartan-moyal-representation}
  Put
  \[
    S_N(f)=\sum_{i=1}^N
      \operatorname{Op}_{\mathrm W}
      (f)(\lambda_i,-\hbar_{\mathrm{W}}\partial_{\lambda_i}).
  \]
  Then, for all \(f,g\in\C[z_1,z_2][[\hbar_{\mathrm{W}}]]\),
  \[
    \hbar_{\mathrm{W}}^{-1}[S_N(f),S_N(g)]=S_N(\{f,g\}_{\hbar_{\mathrm{W}}}).
  \]
\end{lem}

\begin{proof}
  The \(i\)-th summand of \(S_N(f)\) lies in the one-particle Weyl
  algebra generated by \(\lambda_i\) and
  \(-\hbar_{\mathrm{W}}\partial_{\lambda_i}\).  Weyl quantization is an algebra
  isomorphism from the formal Moyal algebra of the constant Poisson
  plane to this one-particle Weyl algebra, so
  \[
    \hbar_{\mathrm{W}}^{-1}[
      \operatorname{Op}_{\mathrm W}(f)_i,
      \operatorname{Op}_{\mathrm W}(g)_i]
    =
    \operatorname{Op}_{\mathrm W}(\{f,g\}_{\hbar_{\mathrm{W}}})_i .
  \]
  Summands with different \(i\) commute, because they use disjoint
  Cartan variables.  Summing over \(i\) gives the displayed identity.
\end{proof}

\begin{thm}[Finite-\(N\) radial-parts identity]
\label{thm:app-radial-finite-N}
  Under the radial-parts datum of
  Statement~\ref{stmt:app-radial-parts-hypotheses}, put
  \[
    S_N(f)=\sum_{i=1}^N
      \operatorname{Op}_{\mathrm W}
      (f)(\lambda_i,-\hbar_{\mathrm{W}}\partial_{\lambda_i}).
  \]
  Then
  \[
    \operatorname{Rad}_0(J_N(f))=\Delta^{-1}S_N(f)\Delta,
  \]
  and for all \(f,g\)
  \[
    D_N(f,g)=
      \hbar_{\mathrm{W}}^{-1}[J_N(f),J_N(g)]-J_N(\{f,g\}_{\hbar_{\mathrm{W}}})
  \]
  lies in the quantum moment-map ideal.  Equivalently, in the
  Vandermonde gauge,
  \[
    \operatorname{Rad}_0(D_N(f,g))
      =
      \hbar_{\mathrm{W}}^{-1}[R_N(f),R_N(g)]-R_N(\{f,g\}_{\hbar_{\mathrm{W}}})
      =0 .
  \]
  Thus \(D_N(f,g)\) vanishes in degree-zero quantum Hamiltonian
  reduction.
\end{thm}

\begin{proof}
  The first displayed identity is the exact image normalization
  included in Statement~\ref{stmt:app-radial-parts-hypotheses}.  The
  same radial-parts datum makes the conjugated radial map injective
  on the reduced invariant quotient, so the commutator defect can be
  tested on the Cartan.
  Lemma~\ref{lem:app-radial-cartan-moyal-representation} gives
  \(\hbar_{\mathrm{W}}^{-1}[S_N(f),S_N(g)]=S_N(\{f,g\}_{\hbar_{\mathrm{W}}})\), and
  Lemma~\ref{lem:app-radial-dunkl-calogero-cancellation} gives the
  same identity after Vandermonde conjugation:
  \[
    \hbar_{\mathrm{W}}^{-1}[R_N(f),R_N(g)]=R_N(\{f,g\}_{\hbar_{\mathrm{W}}}).
  \]
  Thus \(\operatorname{Rad}_0(D_N(f,g))=0\).  The radial-map
  injectivity of Lemma~\ref{lem:app-radial-map-injectivity} places
  \(D_N(f,g)\) in the quantum moment-map ideal.

  Lemma~\ref{lem:app-radial-nonlinear-shear-obstruction} marks the
  exact theorem behind the full displayed formula: mixed Weyl
  images in this convention require the quantum correction beyond
  classical nonlinear shear covariance.  The all-bidegree decorated PBW Stokes
  closure \(\Omega^{\mathrm{rad}}_{a,b}=0\) is given uniformly by
  Theorem~\ref{thm:app-radial-bicomplex-acyclic} under the same
  statement datum and by exact computation on the bidegrees of
  Remark~\ref{rmk:app-radial-finite-verification} by
  Proposition~\ref{prop:app-radial-pr-unconditional-window}; a
  nonzero obstruction is a normalized signed row
  \((\phi,-\lambda)\in\ker B_{a,b}^*\), structurally excluded under
  Statement~\ref{stmt:app-radial-parts-hypotheses} and excluded by exact
  computation on the verified bidegrees.
\end{proof}

\begin{defn}[Stable primitive extraction]
\label{def:app-radial-stable-primitive-extraction}
  Fix a total polynomial degree bound \(D_{\mathrm{poly}}\) and take
  \(N\geq D_{\mathrm{poly}}\).  Let
  \[
    \kappa_{N,D_{\mathrm{poly}}}
  \]
  denote the connected single-trace primitive/indecomposable projection
  on the stable trace algebra in degrees \(\leq D_{\mathrm{poly}}\), for
  the block-sum coproduct \(\Delta_\oplus\): it kills the empty trace and
  all decomposable multi-trace products, equivalently passes to the
  augmentation-ideal quotient \(I/I^2\), and it sends the
  Capelli-renormalized generator
  \[
    J_N(z_1^az_2^b),\qquad 0<a+b\leq D_{\mathrm{poly}},
  \]
  to its single-trace indecomposable class.  The stable map
  \(\kappa_{\infty,D_{\mathrm{poly}}}\) is the stable-tail identification of these
  projections in the Procesi--Razmyslov stable range, after passing
  through the renormalized associated-graded \(J_N\)-basis transition
  system.
\end{defn}

\begin{lem}[Connected-trace projection after Capelli triangularization]
\label{lem:app-radial-connected-projection}
  Fix a polynomial degree bound \(D_{\mathrm{poly}}\).  In the stable range
  \(N\geq D_{\mathrm{poly}}\), the Capelli triangular system
  \[
    J_N(z_1^az_2^b)\equiv
      \sum_{r=0}^{\min(a,b)}
        \left(-\frac{\hbar_{\mathrm{W}} N}{2}\right)^r
        r!\binom ar\binom br\,T_{a-r,b-r}
  \]
  identifies the span of nonempty Weyl traces \(J_N(z_1^az_2^b)\),
  \(a+b\leq D_{\mathrm{poly}}\), with the span of nonempty normal-ordered traces
  \(T_{a,b}\), modulo the empty trace and decomposable multi-traces.
  The primitive/indecomposable projection
  \[
    \kappa_{N,D_{\mathrm{poly}}}\bigl(J_N(z_1^az_2^b)\bigr)
      \longmapsto z_1^az_2^b
  \]
  is therefore degreewise injective on the scalar-reduced connected
  single-trace quotient in the stable range.
\end{lem}

\begin{proof}
  The displayed Capelli formula is triangular for total degree:
  its \(r=0\) term is \(T_{a,b}\), and every \(r\geq1\) term has total
  degree \(a+b-2r\).  The diagonal coefficient is \(1\), hence the
  triangular matrix is unipotent and invertible over
  \(\C[[\hbar_{\mathrm{W}} N]]\).  The only term which can become the empty trace is
  the term with \(a-r=b-r=0\), and the scalar-reduced connected quotient
  removes it.  Decomposable products are removed by the connected
  projection.  Stable Procesi--Razmyslov trace independence in degree
  \(D_{\mathrm{poly}}\) gives linear independence of the remaining connected trace
  generators for \(N\geq D_{\mathrm{poly}}\).  Thus the stated projection is injective
  degreewise in the stable range.
\end{proof}

\begin{thm}[Connected-trace bracket compatibility and stable primitive projections]
\label{thm:app-radial-stable-connected}
  Under the same radial-parts datum as
  Theorem~\ref{thm:app-radial-finite-N}, for each fixed polynomial
  degree the classes
  \[
    \overline{J_N(f)}
  \]
  have a stable connected trace limit after removing the empty
  trace and disconnected products.  This limit is the renormalized
  associated-graded connected transition system associated to the
  finite-rank normal-ordered trace algebras over
  \(\C[[\hbar_{\mathrm{W}}]]\).  In that limit,
  \[
    \hbar_{\mathrm{W}}^{-1}[\partial_{\mathrm{bb},\hbar_{\mathrm{W}}}(f),
      \partial_{\mathrm{bb},\hbar_{\mathrm{W}}}(g)]_{\mathrm{conn,raw}}
      =
      \partial_{\mathrm{bb},\hbar_{\mathrm{W}}}(\{f,g\}_{\hbar_{\mathrm{W}}}).
  \]
  More explicitly, for \(N\geq D_{\mathrm{poly}}\) and all \(f,g\) with
  all monomials in \(f,g,\{f,g\}_{\hbar_{\mathrm{W}}}\) of total degree
  \(\leq D_{\mathrm{poly}}\),
  \[
    \kappa_{N,D_{\mathrm{poly}}}\!\left(
      \hbar_{\mathrm{W}}^{-1}[J_N(f),J_N(g)]
    \right)
    =
    \kappa_{N,D_{\mathrm{poly}}}\!\left(
      J_N(\{f,g\}_{\hbar_{\mathrm{W}}})
    \right),
  \]
  and these identities are compatible under \(N\mapsto N+1\).
\end{thm}

\begin{proof}
  Lemma~\ref{lem:app-radial-connected-projection} gives the
  degreewise stable-rank connected single-trace projection and its
  injectivity after removing the empty trace and decomposable products.
  The finite-\(N\) identity of
  Theorem~\ref{thm:app-radial-finite-N} gives equality in the reduced
  quantum Hamiltonian quotient.  Applying \(\kappa_{N,D_{\mathrm{poly}}}\) gives the
  displayed connected-trace identity, since the primitive projection
  kills only the empty trace and decomposable products and is injective
  on the Capelli-renormalized single-trace span by
  Lemma~\ref{lem:app-radial-connected-projection}.  The identity is
  compatible with adjoining extra Cartan variables, because
  \(S_{N+k}(f)\) is obtained from \(S_N(f)\) by adding commuting
  one-particle summands.  Therefore the bracket identity stabilizes
  degreewise in \(N\).
\end{proof}

\begin{defn}[Reduced open-line Wick graph data]
\label{def:app-reduced-open-line-wick-graph-data}
  In the reduced first-order open-line model with action
  \[
    S_0=\int_\R\operatorname{Tr}(\phi_1\,d\phi_2),
  \]
  use Weyl-ordered boundary vertices \(f(\phi_1,\phi_2)\), and choose
  the zero-mode-removed midpoint propagator
  \[
    G(t,s)=\frac12\operatorname{sgn}(t-s),\qquad
    \langle\phi_1(t)\phi_2(s)\rangle=\hbar_{\mathrm{W}} G(t,s),\qquad
    \langle\phi_2(t)\phi_1(s)\rangle=-\hbar_{\mathrm{W}} G(t,s).
  \]
  Self-contractions at a Weyl vertex are set to zero.  The only
  connected two-vertex graphs are \(r\)-fold cross-contractions between
  the two boundary vertices.
\end{defn}

\begin{thm}[Reduced open-line Wick graph expansion]
\label{thm:app-reduced-open-line-moyal-product}
  With the graph data of
  Definition~\ref{def:app-reduced-open-line-wick-graph-data}, the
  ordered local product for two boundary Hamiltonians at \(t>s\) is
  the formal Moyal product
  \[
    F\star G=m\circ\exp\left(\frac{\hbar_{\mathrm{W}}}{2}P\right)(F\otimes G).
  \]
  More precisely, the graph with exactly \(r\) cross-edges contributes
  \[
    W_r^{>}(f,g)
      =
    \frac{1}{r!}\left(\frac{\hbar_{\mathrm{W}}}{2}\right)^rP^r(f,g),
  \]
  while the reversed ordering \(s>t\) contributes
  \[
    W_r^{<}(f,g)
      =
    \frac{1}{r!}\left(-\frac{\hbar_{\mathrm{W}}}{2}\right)^rP^r(f,g).
  \]
  Hence the antisymmetrized factorization commutator is
  \[
    [f,g]_{\mathrm{graph}}
      =
    \sum_{r\geq0}
      \frac{1-(-1)^r}{2^r r!}\hbar_{\mathrm{W}}^rP^r(f,g)
      =
    \sum_{\substack{r\geq1\\r\ \mathrm{odd}}}
      \frac{2}{2^r r!}\hbar_{\mathrm{W}}^rP^r(f,g).
  \]
  The one-edge and three-edge commutator coefficients are exactly
  \[
    \hbar_{\mathrm{W}} P(f,g),\qquad \frac{\hbar_{\mathrm{W}}^3}{24}P^3(f,g).
  \]
  All even commutator orders vanish.
\end{thm}

\begin{proof}
  For ordered insertions with \(t>s\), the midpoint propagator gives
  \(G(t,s)=1/2\).  A single cross-contraction therefore contributes
  \[
    \frac{\hbar_{\mathrm{W}}}{2}
    \left(\partial_{z_1}\otimes\partial_{z_2}
      -\partial_{z_2}\otimes\partial_{z_1}\right)
    =\frac{\hbar_{\mathrm{W}}}{2}P.
  \]
  Wick expansion of the free first-order model exponentiates the
  cross-contraction operator.  Weyl vertices have zero
  self-contractions; hence a connected two-vertex graph is determined
  by its number \(r\) of cross-edges.  The exponential gives the
  symmetry factor \(1/r!\), and the ordered propagator gives
  \((\hbar_{\mathrm{W}}/2)^r\).  This proves the formula for \(W_r^{>}\).  Reversing
  the order changes the propagator from \(1/2\) to \(-1/2\), giving
  \(W_r^{<}\).  Subtracting the two orderings gives the displayed
  commutator series.  The coefficients in orders \(1\) and \(3\) are
  \(2/(2\cdot1!)=1\) and \(2/(2^3\cdot3!)=1/24\).
\end{proof}

\begin{cor}[All-order reduced Wick realization of Moyal]
\label{cor:app-reduced-graph-realizes-moyal}
  The reduced open-line Wick graph expansion of
  Theorem~\ref{thm:app-reduced-open-line-moyal-product} realizes the
  Weyl/Moyal Lie bracket of
  Definition~\ref{def:app-radial-moyal-convention} to all orders:
  \[
    [f,g]_{\mathrm{graph}}=[f,g]_{\mathrm{raw}},\qquad
    \hbar_{\mathrm{W}}^{-1}[f,g]_{\mathrm{graph}}=\{f,g\}_{\hbar_{\mathrm{W}}} .
  \]
\end{cor}

\begin{proof}
  Compare the graph commutator series in
  Theorem~\ref{thm:app-reduced-open-line-moyal-product} with the formal
  Weyl commutator series in
  Theorem~\ref{thm:app-radial-moyal-coefficients}.  The coefficient of
  \(P^r\) is \(2\hbar_{\mathrm{W}}^r/(2^r r!)\) for odd \(r\) and \(0\) for even
  \(r\) in both series.
\end{proof}

\begin{rmk}[Verified and open bidegree windows]
\label{rmk:app-radial-finite-verification}
  Table~\ref{tab:app-radial-finite-window-status} is the finite-window
  ledger used in the exact rational reduction.  The proved rows are the
  Procesi--Razmyslov kernel-correction route of
  Proposition~\ref{prop:app-radial-pr-unconditional-window}.  The first
  open exact finite systems are \((6,8),(8,6)\) and the balanced case
  \((7,7)\).  This is not a non-vanishing theorem: no signed row in
  \(\ker B_{a,b}^*\) is known for these bidegrees.
  The all-bidegree closure is
  Theorem~\ref{thm:app-radial-bicomplex-acyclic} only under
  Statement~\ref{stmt:app-radial-parts-hypotheses}; without that
  statement, each unlisted bidegree is governed by its exact cokernel
  class \(\Omega^{\mathrm{rad}}_{a,b}\).
\end{rmk}

\subsection*{Exact rational reduction}\label{ssec:app-radial-rational-reduction}

\begin{ex}[The first nonzero balanced cyclic-pivot residual]\label{ex:app-radial-first-balanced-residual}
  At \((a,b)=(2,2)\) and \((3,3)\) the classical lift already lies in
  \(\operatorname{im} D^\square_{a,b}\), and
  \(\Omega^{\mathrm{rad}}_{a,b}=0\) with zero correction.  The first
  bidegree where the cyclic pivot produces a residual is \((4,4)\),
  with residual
  \[
    \tfrac{43}{14}
      \Bigl(
        \overline{\operatorname{Tr}(XXYXYY)}
        -\overline{\operatorname{Tr}(XYXXYY)}
        -\overline{\operatorname{Tr}(XYYXYX)}
        +\overline{\operatorname{Tr}(XYYYXX)}
      \Bigr)
  \]
  killed by a four-term correction in \(\ker\partial\).
\end{ex}

\begin{thm}[Per-bidegree exact rational reduction]\label{thm:app-radial-bidegree-rational-reduction}
  Each bidegree datum is the triple
  \[
    \bigl(\,
        \beta_{a,b},\;
        B_{a,b}\in M_{r_0\times r_1}(\Q),\;
        v_{a,b}\in\Q^{r_0}
      \,\bigr),
  \]
  where \(\beta_{a,b}\) is the indexed-diagram basis ordering,
  \(r_1=\dim\mathsf W_{a,b}=\binom{a+b-2}{a-1}\),
  \(r_0=\dim\bigl(\mathsf W_{a,b}\oplus\mathsf S_{a,b}\bigr)\),
  \(B_{a,b}\) is the Stokes boundary in the basis \(\beta_{a,b}\), and
  \(v_{a,b}=(T_{a,b},E^+_{a,b})\) is the radial comparison vector.
  Bareiss elimination over \(\Q\) decides whether
  \(v_{a,b}\in\operatorname{im}B_{a,b}\) in exact rational
  arithmetic and yields the dimensions
  \[
    \operatorname{rank}B_{a,b},\quad
    \dim\ker B_{a,b}^*,\quad
    \dim\operatorname{coker}B_{a,b}.
  \]
  When \(v_{a,b}\in\operatorname{im}B_{a,b}\), the elimination returns
  a rational correction
  \(C^+_{a,b}\in\Q^{r_1}\) with
  \(B_{a,b}C^+_{a,b}=v_{a,b}\), giving
  \(\Omega^{\mathrm{rad}}_{a,b}=0\) at that bidegree.  In the obstruction
  case, the elimination returns the obstruction covector
  \((\phi,-\lambda)\in\ker B_{a,b}^*\) with the strict pairing
  \(\lambda(E^+_{a,b})-\phi(T_{a,b})\ne0\); the obstruction class
  \(\Omega^{\mathrm{rad}}_{a,b}\in\operatorname{coker}B_{a,b}\) is its
  equivalence class.  Every arithmetic step is exact rational; the
  same elimination on the same basis ordering returns the same
  matrix entries.
\end{thm}

\begin{cor}[Vanishing on the proved finite-window range]\label{cor:app-radial-vanishing-on-listed-range}
  For every bidegree \((a,b)\) in a row marked \emph{proved} in
  Remark~\ref{rmk:app-radial-finite-verification},
  the rational reduction of
  Theorem~\ref{thm:app-radial-bidegree-rational-reduction} gives
  \(v_{a,b}\in\operatorname{im}B_{a,b}\): the comparison vector lies in
  the image of the Stokes boundary, with explicit correction
  \(C^+_{a,b}\) and dimensions
  \(\dim\ker B_{a,b}^*,\,\dim\operatorname{coker}B_{a,b}\) read off from
  the elimination.  In particular
  \(\Omega^{\mathrm{rad}}_{a,b}=0\) on the proved finite-window
  bidegrees.  For a bidegree in the undecided or unlisted rows of
  Remark~\ref{rmk:app-radial-finite-verification}, the per-bidegree class
  \(\Omega^{\mathrm{rad}}_{a,b}\in\operatorname{coker}B_{a,b}\) is the
  exact obstruction whose vanishing is the all-bidegree statement.
\end{cor}

\begin{rmk}[Scalar quotient and the \(U(1)\) anomaly]
\label{rmk:app-radial-scalar-quotient}
  Constants are central in the Moyal Lie algebra and are killed in
  \(\bar A_{\hbar_{\mathrm{W}}}\).  Before this scalar reduction the linear
  commutator carries the central term
  \([z_1,z_2]_*= \hbar_{\mathrm{W}}\), and on the brane side the corresponding
  matrix trace is the Capelli shift \(\hbar_{\mathrm{W}} N\).  The reduced theorem
  is the theorem after quotienting by this scalar line.
  The unreduced line is the \(U(1)\) anomaly class of
  Theorem~\ref{thm:u1-center-anomaly}.
\end{rmk}

\begin{prop}[Moyal--Capelli normalization and realization boundary]
\label{prop:app-moyal-capelli-realization-boundary}
  The following statements are used by the
  degree-zero Moyal sector.
  \begin{enumerate}
    \item The formal Weyl algebra on the constant symplectic plane gives
      the raw commutator coefficients
      \[
        [f,g]_{\mathrm{raw}}^{(2s+1)}
        =
        \frac{\hbar_{\mathrm{W}}^{2s+1}}{2^{2s}(2s+1)!}P^{2s+1}(f,g),
        \qquad
        [f,g]_{\mathrm{raw}}^{(2s)}=0.
      \]
      In particular the third raw coefficient is
      \(\hbar_{\mathrm{W}}^3P^3(f,g)/24\).
    \item The finite-rank brane algebra is the degree-zero normalizer
      quotient of the matrix Weyl algebra by the traceless quantum
      moment ideal.  With the normal-ordering convention of
      Definition~\ref{def:app-radial-matrix-weyl}, the traceless
      comoment congruence of
      Remark~\ref{rmk:app-weyl-versus-matrix-equation} is
      \[
        YX-XY+\hbar_{\mathrm{W}} N I_N\equiv0
        \pmod{\mathcal W_N\widehat\mu(\lie{sl}_N)}.
      \]
      This scalar Capelli shift is represented by the triangular change of
      basis \(T_{a,b}\leftrightarrow J_N(z_1^az_2^b)\), and the remaining
      scalar Hamiltonian line is removed only after passing to
      \(\bar A_{\hbar_{\mathrm{W}}}\) and to the empty-trace quotient.
    \item The equality
      \[
        \hbar_{\mathrm{W}}^{-1}[J_N(f),J_N(g)]
        =
        J_N(\{f,g\}_{\hbar_{\mathrm{W}}})
      \]
      in degree-zero quantum Hamiltonian reduction holds in every
      bidegree under
      Statement~\ref{stmt:app-radial-parts-hypotheses} by
      Theorem~\ref{thm:app-radial-bicomplex-acyclic}, and holds by
      exact computation on the bidegrees listed in
      Remark~\ref{rmk:app-radial-finite-verification} by
      Proposition~\ref{prop:app-radial-pr-unconditional-window}.  The
      obstruction-theoretic equivalence
      \[
        \Omega^{\mathrm{rad}}_{a,b}=0
        \quad\Longleftrightarrow\quad
        D^\square_{a,b}H_{a,b}=R^{\mathrm{free}}_{a,b}
      \]
      holds in every bidegree by
      Proposition~\ref{prop:app-radial-dual-potential-obstruction}.
      The moment-ideal primitive
      \(E_{a,b,N}=\sum A^j{}_i(a,b,N)\widehat\mu^i{}_j\) of
      Lemma~\ref{lem:app-quantum-shear-primitive-obstruction} and the
      class \(\mathfrak o_{a,b}\) of
      Definition~\ref{def:app-free-kernel-correction-obstruction} are
      equivalent intermediate forms after choosing a classical lift.
    \item The quantum principal-part target is the continuous
      \(\bar A_{\hbar_{\mathrm{W}}}\)-coadjoint module.  It is defined by
      \[
        (\operatorname{ad}^{\vee}_{f,\hbar_{\mathrm{W}}}\rho)(g)
        =
        -\rho(\{f,g\}_{\hbar_{\mathrm{W}}}).
      \]
      For \(f=z_1^pz_2^q\) and the principal-part basis
      \(\rho_{a,b}\), this gives the finite formula
      \[
        z_1^pz_2^q\cdot_{\hbar_{\mathrm{W}}}\rho_{a,b}
        =
        -\sum_{\substack{r\geq1\\ r\ \mathrm{odd}}}
        \frac{\hbar_{\mathrm{W}}^{r-1}}{2^{r-1}r!}\,
        C_r(p,q;a-p+r,b-q+r)\,
        \rho_{a-p+r,b-q+r},
      \]
      where
      \[
        C_r(p,q;m,n)=
        \sum_{\alpha=0}^{r}(-1)^\alpha\binom r\alpha
        (p)_{r-\alpha}(q)_\alpha(m)_\alpha(n)_{r-\alpha},
      \]
      and terms with negative indices or scalar target \((0,0)\) are
      zero.  The \(r=1\) part is the classical Matlis coadjoint action.
      The \(r\ge3\) terms are genuine Moyal corrections; the polynomial
      one-\(\psi\) PBW action records only the \(r=1\) term.
    \item The reduced open-line Wick graph computation realizes the same
      formal Moyal product with midpoint propagator.  An unreduced
      Costello graph/QME theorem is the datum consisting of the brane-defect
      propagator, counterterms, anomaly cancellation, and
      \(d_{\mathrm{QME}}\)-homotopies in the local BV complex; those data
      belong to the unreduced BV target.
  \end{enumerate}
\end{prop}

\begin{proof}
  The five items collect already-proven content.  The raw Moyal
  commutator coefficients are
  Theorem~\ref{thm:app-radial-moyal-coefficients}, with the
  \(1/24\) and \(1/1920\) check at
  Example~\ref{ex:app-radial-higher-moyal-check}.  The Capelli
  triangular contact relation is the normal-ordered comoment
  calculation of
  Definition~\ref{def:app-radial-matrix-weyl} and
  Lemma~\ref{lem:app-radial-capelli-triangular}, with the scalar
  quotient of Definition~\ref{def:app-radial-moyal-convention} and
  Remark~\ref{rmk:app-radial-scalar-quotient}.  The Cartan
  one-particle commutator and the all-bidegree compatibility are
  Theorem~\ref{thm:app-radial-finite-N} and
  Theorem~\ref{thm:app-radial-bicomplex-acyclic}; the obstruction-
  theoretic equivalence is
  Proposition~\ref{prop:app-radial-dual-potential-obstruction} with
  primitive form
  Lemma~\ref{lem:app-quantum-shear-primitive-obstruction} and
  fixed-lift Hodge presentation
  Statement~\ref{stmt:app-free-kernel-homotopy-obstruction}; the
  edge strips are
  Proposition~\ref{prop:app-edge-pbw-telescoping}, the exact finite
  window
  Proposition~\ref{prop:app-radial-pr-unconditional-window}, with
  interior anchors
  Propositions~\ref{prop:app-free-trace-diagram-K44},
  \ref{prop:app-free-residual-vanishings}.  The residue-dual coadjoint
  formula in (4) uses the same \(P^r\)-coefficients as (1); falling
  factorials vanish for \(r>p+q\), and the scalar target is removed by
  the Hamiltonian quotient (special fibre
  Proposition~\ref{prop:principal-part-coadjoint}; quantum target
  Construction~\ref{constr:quantum-principal-part-descendant-target}).
  The open-line midpoint Moyal realisation is
  Theorem~\ref{thm:app-reduced-open-line-moyal-product} with
  Corollary~\ref{cor:app-reduced-graph-realizes-moyal}; the separation
  from the Costello graph/QME realization belongs to
  Theorem~\ref{thm:phi-hbar-all-order}.
\end{proof}

\begin{rmk}[Frontier of the all-bidegree obstruction]
\label{rmk:app-radial-open-bidegree-boundary}
  Theorem~\ref{thm:app-radial-bicomplex-acyclic} gives
  \(\Omega^{\mathrm{rad}}_{a,b}=0\) for every interior bidegree
  \((a,b)\) under
  Statement~\ref{stmt:app-radial-parts-hypotheses}; clauses~(1)
  and~(2) are the Levasseur--Stafford radial-kernel theorem, and
  clause~(3) enters only through the Moyal-bracket polynomial
  \(\{z_1^{a+1}/(a+1),z_2^{b+1}\}_{\hbar_{\mathrm{W}}}\) of the chosen bidegree.
  Corollary~\ref{cor:app-radial-bicomplex-clause3-equivalence}
  refines this to a logical equivalence: under clauses~(1)
  and~(2), \(\Omega^{\mathrm{rad}}_{a,b}=0\) is equivalent to
  clause~(3) modulo the moment ideal for that Moyal-bracket
  polynomial.  The bidegree-local obstruction is concentrated in
  clause~(3) for the Moyal-bracket polynomial of the chosen
  bidegree.

  Independently of clause~(3),
  Proposition~\ref{prop:app-radial-pr-unconditional-window}
  closes \(\Omega^{\mathrm{rad}}_{a,b}=0\) by exact rational
  linear algebra against the moment ideal, independently of the
  statement datum.  The bidegrees enumerated in
  Remark~\ref{rmk:app-radial-finite-verification} (edges, all
  non-edge bidegrees of total degree at most \(13\), and the six
  total-degree-\(14\) bidegrees
  \((3,11),(11,3),(4,10),(10,4),(5,9),(9,5)\)) are closed by this
  exact rational normal-form route.  The first undecided exact
  finite systems are \((6,8),(8,6),(7,7)\).

  In the bidegrees verified by the exact finite-window route, the
  signed-row obstruction set of
  Proposition~\ref{prop:app-radial-dual-potential-obstruction}
  is empty; in the remaining bidegrees, its vanishing is
  equivalent to
  Statement~\ref{stmt:app-radial-parts-hypotheses}~clause~(3) for
  the corresponding Moyal-bracket polynomial.  A defect in that
  clause produces a signed dual row
  \((\phi,-\lambda)\in\ker B_{a,b}^*\) with
  \(\lambda(E^+_{a,b})-\phi(T_{a,b})\neq0\) identical with the
  clause-(3) defect inside \(\mathsf{Diag}_{a,b}^{*}\).  An explicit
  normal-form formula realizing clause~(3) for every mixed monomial of
  degree \(\leq a+b-2\) in the Capelli/Weyl convention is the remaining
  all-bidegree radial-image datum; a failure of that datum is exactly
  the signed-row obstruction above.  Levasseur--Stafford spherical-Hecke
  methods provide the Harish--Chandra kernel input, not this
  Capelli/Weyl normalization automatically.

  Reflection symmetry \((a,b)\leftrightarrow(b,a)\) is implemented
  by the symplectic involution \(z_1\leftrightarrow z_2\),
  \(X\leftrightarrow Y\), with sign
  \(\omega(z_1,z_2)\to-\omega(z_2,z_1)\) recorded in
  Proposition~\ref{prop:app-local-geometry-scalar-shadow}, and is
  preserved by the radial-bicomplex argument.
\end{rmk}
